% !TEX program = lualatex
% ---------------------------------------------------------------------------------------------
%  Readback note 24: The characters, the level prime and the index sets of the data; the grid spaces
%
%  Build:  python3 tools/rb.py note <slug>
%
%  One unit of a Lean artifact, rewritten as ordinary mathematics. The prose is hand-written. Lean names,
%  line ranges and counts come from data.tex; quoted Lean statements are copied from the artifact by the build.
%
%  Rules of the text: standard terminology; every symbol is explained where it first appears and listed again
%  in the notation at the end; no letter has two meanings; a proof is complete or says what it takes for granted;
%  Lean names appear only in the LEAN lines and in the concordance.
%
%  The signature rule (the build enforces it): a function, relation or predicate that the artifact defines is
%  written only through its notation macro, which the build makes from the spec (`notation`) with one slot for
%  each explicit argument of the Lean definition. No argument is fixed and dropped, also not "for brevity" and
%  not for one section. Every argument and every hypothesis of a stated declaration stands inside its
%  statement, or in a Setting that the statement cites by number; only implicit and instance arguments (the
%  ambient objects) may stand in the conventions.
% ---------------------------------------------------------------------------------------------
\documentclass[11pt,a4paper]{article}
\ifdefined\pdfvariable\pdfvariable suppressoptionalinfo \numexpr 1+2+4+8\relax\fi
\usepackage[a4paper,hmargin=35mm,top=22mm,bottom=27mm,footskip=12mm,nohead]{geometry}

\usepackage{fontspec}
\directlua{
  luaotfload.add_fallback("gfallback", {"STIXTwoMath-Regular.otf:mode=node;", "DejaVuSans.ttf:mode=node;"})
  luaotfload.add_fallback("gmonofallback", {"DejaVuSansMono.ttf:mode=node;", "STIXTwoMath-Regular.otf:mode=node;"})
}
\setmainfont{NewCM10}[Extension=.otf,UprightFont=*-Book,ItalicFont=*-BookItalic,
  BoldFont=*-Bold,BoldItalicFont=*-BoldItalic,SmallCapsFont=*-Book,SmallCapsFeatures={Letters=SmallCaps},
  RawFeature={fallback=gfallback}]
\setsansfont{NewCMSans10}[Extension=.otf,UprightFont=*-Book,ItalicFont=*-BookOblique,
  BoldFont=*-Bold,BoldItalicFont=*-BoldOblique,RawFeature={fallback=gfallback}]
\setmonofont{JetBrainsMono}[Extension=.otf,UprightFont=*-Regular,ItalicFont=*-Italic,
  BoldFont=*-SemiBold,BoldItalicFont=*-SemiBoldItalic,Scale=MatchLowercase,
  RawFeature={-calt;-liga;fallback=gmonofallback}]
\usepackage{amsmath}
\usepackage{amsthm}
\usepackage[math-style=TeX]{unicode-math}
\setmathfont{NewCMMath-Book.otf}
\emergencystretch=2.5em\relax

\usepackage{xcolor}
\usepackage{tikz}
\usetikzlibrary{calc,arrows.meta,positioning,backgrounds}
\usepackage[protrusion=true,expansion=true]{microtype}
\usepackage{enumitem}
\usepackage{needspace}
\usepackage{booktabs}
\usepackage{array}
\usepackage{longtable}
\usepackage{fancyhdr}
\usepackage{lastpage}
\usepackage{multicol}
\usepackage[hidelinks,pdfusetitle]{hyperref}

\definecolor{ink}{HTML}{0F172A}
\definecolor{inksoft}{HTML}{475569}
\definecolor{inkfaint}{HTML}{94A3B8}
\definecolor{rulegrey}{HTML}{CBD5E1}
\definecolor{arith}{HTML}{C2410C}
\definecolor{est}{HTML}{7C3AED}
\definecolor{ana}{HTML}{0891B2}
\definecolor{ok}{HTML}{047857}
\definecolor{leancol}{HTML}{3730A3}

\title{The characters, the level prime and the index sets of the data; the grid spaces}
\author{grh-compact}
\date{}

% ------------------------------------------------------------------------------------- macros
\newcommand{\leanfont}{\ttfamily\fontsize{8.6}{10.5}\selectfont\color{leancol}}
\newcommand{\lean}[1]{\csname LN@#1\endcsname}
\newcommand{\fillin}[1]{{\color{arith}\itshape[#1]}}  % a place of the template that is not written yet; the build reports every one that is left
\newcommand{\leanquote}[1]{\csname LQ@#1\endcsname}
% generated by rb.py note; do not edit
\newcommand{\ConcDecls}{88}
\newcommand{\Pin}{72fead10}
\newcommand{\PinText}{that of commit {\leanfont 72fead10}}
\newcommand{\ArtifactName}{grh-compact}
\expandafter\newcommand\csname LN@hasLogBoundvarpi\endcsname{\href{https://github.com/lsernau/grh-lean/blob/72fead102f663d5c7a407964a13dc0cf9d977207/grh-compact/src/New/Seesaw/IndexTypes.lean\#L19-L27}{{\leanfont Data.\allowbreak{}hasLogBound\_varpi}}}
\expandafter\newcommand\csname LN@hasLogBoundrootDisc\endcsname{\href{https://github.com/lsernau/grh-lean/blob/72fead102f663d5c7a407964a13dc0cf9d977207/grh-compact/src/New/Seesaw/IndexTypes.lean\#L30-L33}{{\leanfont Data.\allowbreak{}hasLogBound\_rootDisc}}}
\expandafter\newcommand\csname LN@piGBar\endcsname{\href{https://github.com/lsernau/grh-lean/blob/72fead102f663d5c7a407964a13dc0cf9d977207/grh-compact/src/New/Seesaw/IndexTypes.lean\#L37-L37}{{\leanfont Data.\allowbreak{}piGBar}}}
\expandafter\newcommand\csname LN@IsChosen\endcsname{\href{https://github.com/lsernau/grh-lean/blob/72fead102f663d5c7a407964a13dc0cf9d977207/grh-compact/src/New/Seesaw/IndexTypes.lean\#L41-L41}{{\leanfont Data.\allowbreak{}IsChosen}}}
\expandafter\newcommand\csname LN@ChosenPlace\endcsname{\href{https://github.com/lsernau/grh-lean/blob/72fead102f663d5c7a407964a13dc0cf9d977207/grh-compact/src/New/Seesaw/IndexTypes.lean\#L45-L45}{{\leanfont Data.\allowbreak{}ChosenPlace}}}
\expandafter\newcommand\csname LN@piGBarmulstar\endcsname{\href{https://github.com/lsernau/grh-lean/blob/72fead102f663d5c7a407964a13dc0cf9d977207/grh-compact/src/New/Seesaw/IndexTypes.lean\#L47-L48}{{\leanfont Data.\allowbreak{}piGBar\_mul\_star}}}
\expandafter\newcommand\csname LN@piGnezero\endcsname{\href{https://github.com/lsernau/grh-lean/blob/72fead102f663d5c7a407964a13dc0cf9d977207/grh-compact/src/New/Seesaw/IndexTypes.lean\#L50-L50}{{\leanfont Data.\allowbreak{}piG\_ne\_zero}}}
\expandafter\newcommand\csname LN@piGBarnezero\endcsname{\href{https://github.com/lsernau/grh-lean/blob/72fead102f663d5c7a407964a13dc0cf9d977207/grh-compact/src/New/Seesaw/IndexTypes.lean\#L52-L52}{{\leanfont Data.\allowbreak{}piGBar\_ne\_zero}}}
\expandafter\newcommand\csname LN@existscofactor\endcsname{\href{https://github.com/lsernau/grh-lean/blob/72fead102f663d5c7a407964a13dc0cf9d977207/grh-compact/src/New/Seesaw/IndexTypes.lean\#L54-L55}{{\leanfont Data.\allowbreak{}exists\_cofactor}}}
\expandafter\newcommand\csname LN@isChoseniff\endcsname{\href{https://github.com/lsernau/grh-lean/blob/72fead102f663d5c7a407964a13dc0cf9d977207/grh-compact/src/New/Seesaw/IndexTypes.lean\#L57-L58}{{\leanfont Data.\allowbreak{}isChosen\_iff}}}
\expandafter\newcommand\csname LN@qeqg\endcsname{\href{https://github.com/lsernau/grh-lean/blob/72fead102f663d5c7a407964a13dc0cf9d977207/grh-compact/src/New/Seesaw/IndexTypes.lean\#L62-L63}{{\leanfont Data.\allowbreak{}q\_eq\_g}}}
\expandafter\newcommand\csname LN@finiteisChosen\endcsname{\href{https://github.com/lsernau/grh-lean/blob/72fead102f663d5c7a407964a13dc0cf9d977207/grh-compact/src/New/Seesaw/IndexTypes.lean\#L65-L66}{{\leanfont Data.\allowbreak{}finite\_isChosen}}}
\expandafter\newcommand\csname LN@existsresiduecharoforder\endcsname{\href{https://github.com/lsernau/grh-lean/blob/72fead102f663d5c7a407964a13dc0cf9d977207/grh-compact/src/New/Seesaw/IndexTypes.lean\#L70-L72}{{\leanfont Data.\allowbreak{}exists\_residue\_char\_of\_order}}}
\expandafter\newcommand\csname LN@decla\endcsname{\href{https://github.com/lsernau/grh-lean/blob/72fead102f663d5c7a407964a13dc0cf9d977207/grh-compact/src/New/Seesaw/IndexTypes.lean\#L76-L76}{{\leanfont Data.\allowbreak{}ω}}}
\expandafter\newcommand\csname LN@orderOf\endcsname{\href{https://github.com/lsernau/grh-lean/blob/72fead102f663d5c7a407964a13dc0cf9d977207/grh-compact/src/New/Seesaw/IndexTypes.lean\#L79-L80}{{\leanfont Data.\allowbreak{}orderOf\_ω}}}
\expandafter\newcommand\csname LN@omegaeqoneofnotisChosen\endcsname{\href{https://github.com/lsernau/grh-lean/blob/72fead102f663d5c7a407964a13dc0cf9d977207/grh-compact/src/New/Seesaw/IndexTypes.lean\#L82-L82}{{\leanfont Data.\allowbreak{}omega\_eq\_one\_of\_not\_isChosen}}}
\expandafter\newcommand\csname LN@omegazpowmulpoweqone\endcsname{\href{https://github.com/lsernau/grh-lean/blob/72fead102f663d5c7a407964a13dc0cf9d977207/grh-compact/src/New/Seesaw/IndexTypes.lean\#L84-L85}{{\leanfont Data.\allowbreak{}omega\_zpow\_mul\_pow\_eq\_one}}}
\expandafter\newcommand\csname LN@existscharhasWeight\endcsname{\href{https://github.com/lsernau/grh-lean/blob/72fead102f663d5c7a407964a13dc0cf9d977207/grh-compact/src/New/Seesaw/IndexTypes.lean\#L89-L97}{{\leanfont Data.\allowbreak{}exists\_char\_hasWeight}}}
\expandafter\newcommand\csname LN@muFlat\endcsname{\href{https://github.com/lsernau/grh-lean/blob/72fead102f663d5c7a407964a13dc0cf9d977207/grh-compact/src/New/Seesaw/IndexTypes.lean\#L101-L101}{{\leanfont Data.\allowbreak{}muFlat}}}
\expandafter\newcommand\csname LN@declb\endcsname{\href{https://github.com/lsernau/grh-lean/blob/72fead102f663d5c7a407964a13dc0cf9d977207/grh-compact/src/New/Seesaw/IndexTypes.lean\#L105-L105}{{\leanfont Data.\allowbreak{}χ}}}
\expandafter\newcommand\csname LN@muFlathasWeight\endcsname{\href{https://github.com/lsernau/grh-lean/blob/72fead102f663d5c7a407964a13dc0cf9d977207/grh-compact/src/New/Seesaw/IndexTypes.lean\#L108-L108}{{\leanfont Data.\allowbreak{}muFlat\_hasWeight}}}
\expandafter\newcommand\csname LN@chihasWeight\endcsname{\href{https://github.com/lsernau/grh-lean/blob/72fead102f663d5c7a407964a13dc0cf9d977207/grh-compact/src/New/Seesaw/IndexTypes.lean\#L111-L111}{{\leanfont Data.\allowbreak{}chi\_hasWeight}}}
\expandafter\newcommand\csname LN@Pair\endcsname{\href{https://github.com/lsernau/grh-lean/blob/72fead102f663d5c7a407964a13dc0cf9d977207/grh-compact/src/New/Seesaw/IndexTypes.lean\#L115-L122}{{\leanfont Data.\allowbreak{}Pair}}}
\expandafter\newcommand\csname LN@rowParameter\endcsname{\href{https://github.com/lsernau/grh-lean/blob/72fead102f663d5c7a407964a13dc0cf9d977207/grh-compact/src/New/Seesaw/IndexTypes.lean\#L126-L126}{{\leanfont Data.\allowbreak{}rowParameter}}}
\expandafter\newcommand\csname LN@logLevelTarget\endcsname{\href{https://github.com/lsernau/grh-lean/blob/72fead102f663d5c7a407964a13dc0cf9d977207/grh-compact/src/New/Seesaw/IndexTypes.lean\#L133-L134}{{\leanfont Data.\allowbreak{}logLevelTarget}}}
\expandafter\newcommand\csname LN@logLevelTargetmemIcc\endcsname{\href{https://github.com/lsernau/grh-lean/blob/72fead102f663d5c7a407964a13dc0cf9d977207/grh-compact/src/New/Seesaw/IndexTypes.lean\#L137-L141}{{\leanfont Data.\allowbreak{}logLevelTarget\_mem\_Icc}}}
\expandafter\newcommand\csname LN@N\endcsname{\href{https://github.com/lsernau/grh-lean/blob/72fead102f663d5c7a407964a13dc0cf9d977207/grh-compact/src/New/Seesaw/IndexTypes.lean\#L145-L146}{{\leanfont Data.\allowbreak{}N}}}
\expandafter\newcommand\csname LN@Nprime\endcsname{\href{https://github.com/lsernau/grh-lean/blob/72fead102f663d5c7a407964a13dc0cf9d977207/grh-compact/src/New/Seesaw/IndexTypes.lean\#L149-L149}{{\leanfont Data.\allowbreak{}N\_prime}}}
\expandafter\newcommand\csname LN@NmodEqone\endcsname{\href{https://github.com/lsernau/grh-lean/blob/72fead102f663d5c7a407964a13dc0cf9d977207/grh-compact/src/New/Seesaw/IndexTypes.lean\#L151-L152}{{\leanfont Data.\allowbreak{}N\_modEq\_one}}}
\expandafter\newcommand\csname LN@logNmemwindow\endcsname{\href{https://github.com/lsernau/grh-lean/blob/72fead102f663d5c7a407964a13dc0cf9d977207/grh-compact/src/New/Seesaw/IndexTypes.lean\#L157-L158}{{\leanfont Data.\allowbreak{}log\_N\_mem\_window}}}
\expandafter\newcommand\csname LN@logNmemIcc\endcsname{\href{https://github.com/lsernau/grh-lean/blob/72fead102f663d5c7a407964a13dc0cf9d977207/grh-compact/src/New/Seesaw/IndexTypes.lean\#L161-L162}{{\leanfont Data.\allowbreak{}log\_N\_mem\_Icc}}}
\expandafter\newcommand\csname LN@existspimulstar\endcsname{\href{https://github.com/lsernau/grh-lean/blob/72fead102f663d5c7a407964a13dc0cf9d977207/grh-compact/src/New/Seesaw/IndexTypes.lean\#L173-L175}{{\leanfont Data.\allowbreak{}exists\_pi\_mul\_star}}}
\expandafter\newcommand\csname LN@declc\endcsname{\href{https://github.com/lsernau/grh-lean/blob/72fead102f663d5c7a407964a13dc0cf9d977207/grh-compact/src/New/Seesaw/IndexTypes.lean\#L178-L178}{{\leanfont Data.\allowbreak{}π}}}
\expandafter\newcommand\csname LN@mulstar\endcsname{\href{https://github.com/lsernau/grh-lean/blob/72fead102f663d5c7a407964a13dc0cf9d977207/grh-compact/src/New/Seesaw/IndexTypes.lean\#L182-L182}{{\leanfont Data.\allowbreak{}π\_mul\_star}}}
\expandafter\newcommand\csname LN@TameChar\endcsname{\href{https://github.com/lsernau/grh-lean/blob/72fead102f663d5c7a407964a13dc0cf9d977207/grh-compact/src/New/Seesaw/IndexTypes.lean\#L186-L192}{{\leanfont Data.\allowbreak{}TameChar}}}
\expandafter\newcommand\csname LQ@TameChar\endcsname{\leanquoteline{0}{structure TameChar where}\leanquoteline{2}{toAutChar : AutChar D.E}\leanquoteline{2}{Q : Finset \_}\leanquoteline{2}{(ψ) hasWeight : toAutChar.HasWeight 0 ψ}\leanquoteline{2}{ψ\_eq\_one : ∀ w ∉ Q, ψ w = 1}\leanquoteline{2}{star\_notMem : ∀ w ∈ Q, star w ∉ Q}\leanquoteline{2}{coprime : (D.gProd * D.N).Coprime (∏ w ∈ Q, w.q)}}
\expandafter\newcommand\csname LN@levelNorm\endcsname{\href{https://github.com/lsernau/grh-lean/blob/72fead102f663d5c7a407964a13dc0cf9d977207/grh-compact/src/New/Seesaw/IndexTypes.lean\#L195-L195}{{\leanfont Data.\allowbreak{}TameChar.\allowbreak{}levelNorm}}}
\expandafter\newcommand\csname LN@Twist\endcsname{\href{https://github.com/lsernau/grh-lean/blob/72fead102f663d5c7a407964a13dc0cf9d977207/grh-compact/src/New/Seesaw/IndexTypes.lean\#L199-L199}{{\leanfont Data.\allowbreak{}Twist}}}
\expandafter\newcommand\csname LN@one\endcsname{\href{https://github.com/lsernau/grh-lean/blob/72fead102f663d5c7a407964a13dc0cf9d977207/grh-compact/src/New/Seesaw/IndexTypes.lean\#L206-L206}{{\leanfont Data.\allowbreak{}TameChar.\allowbreak{}one}}}
\expandafter\newcommand\csname LN@decld\endcsname{\href{https://github.com/lsernau/grh-lean/blob/72fead102f663d5c7a407964a13dc0cf9d977207/grh-compact/src/New/Seesaw/IndexTypes.lean\#L210-L210}{{\leanfont Data.\allowbreak{}μ₀}}}
\expandafter\newcommand\csname LQ@decld\endcsname{\leanquoteline{0}{def μ₀ := D.muFlat * (WeilData.upsilonU1 \textasciicircum{} (-D.h : ℤ) * D.χ) \textasciicircum{} j}}
\expandafter\newcommand\csname LN@muhasWeight\endcsname{\href{https://github.com/lsernau/grh-lean/blob/72fead102f663d5c7a407964a13dc0cf9d977207/grh-compact/src/New/Seesaw/IndexTypes.lean\#L213-L215}{{\leanfont Data.\allowbreak{}mu₀\_hasWeight}}}
\expandafter\newcommand\csname LN@liftChar\endcsname{\href{https://github.com/lsernau/grh-lean/blob/72fead102f663d5c7a407964a13dc0cf9d977207/grh-compact/src/New/Seesaw/IndexTypes.lean\#L218-L218}{{\leanfont Data.\allowbreak{}liftChar}}}
\expandafter\newcommand\csname LN@liftTamePart\endcsname{\href{https://github.com/lsernau/grh-lean/blob/72fead102f663d5c7a407964a13dc0cf9d977207/grh-compact/src/New/Seesaw/IndexTypes.lean\#L222-L222}{{\leanfont Data.\allowbreak{}liftTamePart}}}
\expandafter\newcommand\csname LN@hasWeightliftChar\endcsname{\href{https://github.com/lsernau/grh-lean/blob/72fead102f663d5c7a407964a13dc0cf9d977207/grh-compact/src/New/Seesaw/IndexTypes.lean\#L224-L225}{{\leanfont Data.\allowbreak{}hasWeight\_liftChar}}}
\expandafter\newcommand\csname LN@LiftIndex\endcsname{\href{https://github.com/lsernau/grh-lean/blob/72fead102f663d5c7a407964a13dc0cf9d977207/grh-compact/src/New/Seesaw/IndexTypes.lean\#L228-L228}{{\leanfont Data.\allowbreak{}LiftIndex}}}
\expandafter\newcommand\csname LN@BlockIndex\endcsname{\href{https://github.com/lsernau/grh-lean/blob/72fead102f663d5c7a407964a13dc0cf9d977207/grh-compact/src/New/Seesaw/IndexTypes.lean\#L231-L231}{{\leanfont BlockIndex}}}
\expandafter\newcommand\csname LN@firstBlock\endcsname{\href{https://github.com/lsernau/grh-lean/blob/72fead102f663d5c7a407964a13dc0cf9d977207/grh-compact/src/New/Seesaw/IndexTypes.lean\#L234-L235}{{\leanfont BlockIndex.\allowbreak{}firstBlock}}}
\expandafter\newcommand\csname LN@equivSum\endcsname{\href{https://github.com/lsernau/grh-lean/blob/72fead102f663d5c7a407964a13dc0cf9d977207/grh-compact/src/New/Seesaw/IndexTypes.lean\#L239-L241}{{\leanfont BlockIndex.\allowbreak{}equivSum}}}
\expandafter\newcommand\csname LN@instFintypeYT\endcsname{\href{https://github.com/lsernau/grh-lean/blob/72fead102f663d5c7a407964a13dc0cf9d977207/grh-compact/src/New/Seesaw/IndexTypes.lean\#L245-L246}{{\leanfont instFintypeYT}}}
\expandafter\newcommand\csname LN@card\endcsname{\href{https://github.com/lsernau/grh-lean/blob/72fead102f663d5c7a407964a13dc0cf9d977207/grh-compact/src/New/Seesaw/IndexTypes.lean\#L248-L249}{{\leanfont BlockIndex.\allowbreak{}card}}}
\expandafter\newcommand\csname LN@decle\endcsname{\href{https://github.com/lsernau/grh-lean/blob/72fead102f663d5c7a407964a13dc0cf9d977207/grh-compact/src/New/Seesaw/IndexTypes.lean\#L252-L252}{{\leanfont Data.\allowbreak{}ϱ}}}
\expandafter\newcommand\csname LN@InitIndex\endcsname{\href{https://github.com/lsernau/grh-lean/blob/72fead102f663d5c7a407964a13dc0cf9d977207/grh-compact/src/New/Seesaw/IndexTypes.lean\#L255-L255}{{\leanfont Data.\allowbreak{}InitIndex}}}
\expandafter\newcommand\csname LN@X\endcsname{\href{https://github.com/lsernau/grh-lean/blob/72fead102f663d5c7a407964a13dc0cf9d977207/grh-compact/src/New/Seesaw/IndexTypes.lean\#L259-L259}{{\leanfont Data.\allowbreak{}X}}}
\expandafter\newcommand\csname LN@shape\endcsname{\href{https://github.com/lsernau/grh-lean/blob/72fead102f663d5c7a407964a13dc0cf9d977207/grh-compact/src/New/Seesaw/IndexTypes.lean\#L263-L263}{{\leanfont Data.\allowbreak{}shape}}}
\expandafter\newcommand\csname LN@RowIndex\endcsname{\href{https://github.com/lsernau/grh-lean/blob/72fead102f663d5c7a407964a13dc0cf9d977207/grh-compact/src/New/Seesaw/IndexTypes.lean\#L267-L267}{{\leanfont Data.\allowbreak{}RowIndex}}}
\expandafter\newcommand\csname LN@Y\endcsname{\href{https://github.com/lsernau/grh-lean/blob/72fead102f663d5c7a407964a13dc0cf9d977207/grh-compact/src/New/Seesaw/IndexTypes.lean\#L270-L270}{{\leanfont Data.\allowbreak{}Y}}}
\expandafter\newcommand\csname LN@cardrowIndex\endcsname{\href{https://github.com/lsernau/grh-lean/blob/72fead102f663d5c7a407964a13dc0cf9d977207/grh-compact/src/New/Seesaw/IndexTypes.lean\#L274-L274}{{\leanfont Data.\allowbreak{}card\_rowIndex}}}
\expandafter\newcommand\csname LN@basisNumber\endcsname{\href{https://github.com/lsernau/grh-lean/blob/72fead102f663d5c7a407964a13dc0cf9d977207/grh-compact/src/New/Seesaw/IndexTypes.lean\#L281-L281}{{\leanfont Data.\allowbreak{}basisNumber}}}
\expandafter\newcommand\csname LN@paramW\endcsname{\href{https://github.com/lsernau/grh-lean/blob/72fead102f663d5c7a407964a13dc0cf9d977207/grh-compact/src/New/Seesaw/IndexTypes.lean\#L285-L285}{{\leanfont Data.\allowbreak{}paramW}}}
\expandafter\newcommand\csname LN@entrySign\endcsname{\href{https://github.com/lsernau/grh-lean/blob/72fead102f663d5c7a407964a13dc0cf9d977207/grh-compact/src/New/Seesaw/IndexTypes.lean\#L289-L289}{{\leanfont Data.\allowbreak{}entrySign}}}
\expandafter\newcommand\csname LN@iteoddeqsubemod\endcsname{\href{https://github.com/lsernau/grh-lean/blob/72fead102f663d5c7a407964a13dc0cf9d977207/grh-compact/src/New/Seesaw/IndexTypes.lean\#L291-L291}{{\leanfont Data.\allowbreak{}ite\_odd\_eq\_sub\_emod}}}
\expandafter\newcommand\csname LN@signsW\endcsname{\href{https://github.com/lsernau/grh-lean/blob/72fead102f663d5c7a407964a13dc0cf9d977207/grh-compact/src/New/Seesaw/IndexTypes.lean\#L295-L295}{{\leanfont Data.\allowbreak{}signsW}}}
\expandafter\newcommand\csname LN@chosen\endcsname{\href{https://github.com/lsernau/grh-lean/blob/72fead102f663d5c7a407964a13dc0cf9d977207/grh-compact/src/New/Seesaw/IndexTypes.lean\#L299-L299}{{\leanfont Data.\allowbreak{}chosen}}}
\expandafter\newcommand\csname LN@G\endcsname{\href{https://github.com/lsernau/grh-lean/blob/72fead102f663d5c7a407964a13dc0cf9d977207/grh-compact/src/New/Seesaw/IndexTypes.lean\#L302-L302}{{\leanfont Data.\allowbreak{}G}}}
\expandafter\newcommand\csname LN@goodTupleFamilies\endcsname{\href{https://github.com/lsernau/grh-lean/blob/72fead102f663d5c7a407964a13dc0cf9d977207/grh-compact/src/New/Seesaw/IndexTypes.lean\#L308-L310}{{\leanfont Data.\allowbreak{}goodTupleFamilies}}}
\expandafter\newcommand\csname LN@e\endcsname{\href{https://github.com/lsernau/grh-lean/blob/72fead102f663d5c7a407964a13dc0cf9d977207/grh-compact/src/New/Seesaw/IndexTypes.lean\#L314-L314}{{\leanfont Data.\allowbreak{}e₀}}}
\expandafter\newcommand\csname LN@existsbranchingpairing\endcsname{\href{https://github.com/lsernau/grh-lean/blob/72fead102f663d5c7a407964a13dc0cf9d977207/grh-compact/src/New/Seesaw/PairTensorData.lean\#L22-L58}{{\leanfont exists\_branching\_pairing}}}
\expandafter\newcommand\csname LN@gridSpace\endcsname{\href{https://github.com/lsernau/grh-lean/blob/72fead102f663d5c7a407964a13dc0cf9d977207/grh-compact/src/New/Seesaw/PairTensorData.lean\#L68-L71}{{\leanfont Data.\allowbreak{}gridSpace}}}
\expandafter\newcommand\csname LQ@gridSpace\endcsname{\leanquoteline{0}{def gridSpace : ∀ o z i j, Space D.E \_ (D.gridSigns z o i j)}\leanquoteline{2}{| 0, true, 0, 0 => ((D.spaceY 0).sum (signedSpace ..)).reindex (.refl \_) D.gridSigns\_zero.symm}\leanquoteline{2}{| 0, false, 0, j => D.spaceY j | \_ + 1, \_, 0, 0 => (gridSpace ..).neg | \_, \_, \_, \_ =>}\leanquoteline{4}{signedSpace ..}}
\expandafter\newcommand\csname LN@gridKernel\endcsname{\href{https://github.com/lsernau/grh-lean/blob/72fead102f663d5c7a407964a13dc0cf9d977207/grh-compact/src/New/Seesaw/PairTensorData.lean\#L75-L75}{{\leanfont Data.\allowbreak{}gridKernel}}}
\expandafter\newcommand\csname LN@GridForm\endcsname{\href{https://github.com/lsernau/grh-lean/blob/72fead102f663d5c7a407964a13dc0cf9d977207/grh-compact/src/New/Seesaw/PairTensorData.lean\#L79-L79}{{\leanfont Data.\allowbreak{}GridForm}}}
\expandafter\newcommand\csname LN@gridChar\endcsname{\href{https://github.com/lsernau/grh-lean/blob/72fead102f663d5c7a407964a13dc0cf9d977207/grh-compact/src/New/Seesaw/PairTensorData.lean\#L83-L84}{{\leanfont Data.\allowbreak{}gridChar}}}
\expandafter\newcommand\csname LN@gridTamePart\endcsname{\href{https://github.com/lsernau/grh-lean/blob/72fead102f663d5c7a407964a13dc0cf9d977207/grh-compact/src/New/Seesaw/PairTensorData.lean\#L88-L88}{{\leanfont Data.\allowbreak{}gridTamePart}}}
\expandafter\newcommand\csname LN@triangleOrder\endcsname{\href{https://github.com/lsernau/grh-lean/blob/72fead102f663d5c7a407964a13dc0cf9d977207/grh-compact/src/New/Seesaw/PairTensorData.lean\#L92-L92}{{\leanfont Data.\allowbreak{}triangleOrder}}}
\expandafter\newcommand\csname LN@Node\endcsname{\href{https://github.com/lsernau/grh-lean/blob/72fead102f663d5c7a407964a13dc0cf9d977207/grh-compact/src/New/Seesaw/PairTensorData.lean\#L96-L96}{{\leanfont Data.\allowbreak{}Node}}}
\expandafter\newcommand\csname LN@NodeIndex\endcsname{\href{https://github.com/lsernau/grh-lean/blob/72fead102f663d5c7a407964a13dc0cf9d977207/grh-compact/src/New/Seesaw/PairTensorData.lean\#L99-L99}{{\leanfont Data.\allowbreak{}NodeIndex}}}
\expandafter\newcommand\csname LN@nodeLevel\endcsname{\href{https://github.com/lsernau/grh-lean/blob/72fead102f663d5c7a407964a13dc0cf9d977207/grh-compact/src/New/Seesaw/PairTensorData.lean\#L103-L104}{{\leanfont Data.\allowbreak{}nodeLevel}}}
\expandafter\newcommand\csname LN@levelIndex\endcsname{\href{https://github.com/lsernau/grh-lean/blob/72fead102f663d5c7a407964a13dc0cf9d977207/grh-compact/src/New/Seesaw/PairTensorData.lean\#L108-L108}{{\leanfont Data.\allowbreak{}levelIndex}}}
\expandafter\newcommand\csname LN@gridLevelForms\endcsname{\href{https://github.com/lsernau/grh-lean/blob/72fead102f663d5c7a407964a13dc0cf9d977207/grh-compact/src/New/Seesaw/PairTensorData.lean\#L112-L113}{{\leanfont Data.\allowbreak{}gridLevelForms}}}
\expandafter\newcommand\csname LN@isConjKernelPairgridKernel\endcsname{\href{https://github.com/lsernau/grh-lean/blob/72fead102f663d5c7a407964a13dc0cf9d977207/grh-compact/src/New/Seesaw/PairTensorData.lean\#L115-L117}{{\leanfont Data.\allowbreak{}isConjKernelPair\_gridKernel}}}
\expandafter\newcommand\csname LN@gridCharhasWeight\endcsname{\href{https://github.com/lsernau/grh-lean/blob/72fead102f663d5c7a407964a13dc0cf9d977207/grh-compact/src/New/Seesaw/PairTensorData.lean\#L119-L120}{{\leanfont Data.\allowbreak{}gridChar\_hasWeight}}}
\expandafter\newcommand\csname LN@isEdgeFamilygrid\endcsname{\href{https://github.com/lsernau/grh-lean/blob/72fead102f663d5c7a407964a13dc0cf9d977207/grh-compact/src/New/Seesaw/PairTensorData.lean\#L125-L128}{{\leanfont Data.\allowbreak{}isEdgeFamily\_grid}}}
\expandafter\newcommand\csname LN@gridSpacediag\endcsname{\href{https://github.com/lsernau/grh-lean/blob/72fead102f663d5c7a407964a13dc0cf9d977207/grh-compact/src/New/Seesaw/PairTensorData.lean\#L132-L136}{{\leanfont Data.\allowbreak{}gridSpace\_diag}}}
\expandafter\newcommand\csname LN@gridSpacesubdiag\endcsname{\href{https://github.com/lsernau/grh-lean/blob/72fead102f663d5c7a407964a13dc0cf9d977207/grh-compact/src/New/Seesaw/PairTensorData.lean\#L140-L143}{{\leanfont Data.\allowbreak{}gridSpace\_subdiag}}}
\expandafter\newcommand\csname LN@pointDescent\endcsname{\href{https://github.com/lsernau/grh-lean/blob/72fead102f663d5c7a407964a13dc0cf9d977207/grh-compact/src/New/Seesaw/PairTensorData.lean\#L148-L150}{{\leanfont Data.\allowbreak{}pointDescent}}}
\expandafter\newcommand\csname LN@iotaPlus\endcsname{\href{https://github.com/lsernau/grh-lean/blob/72fead102f663d5c7a407964a13dc0cf9d977207/grh-compact/src/New/Seesaw/PairTensorData.lean\#L154-L155}{{\leanfont Data.\allowbreak{}iotaPlus}}}
\expandafter\newcommand\csname LN@descentMove\endcsname{\href{https://github.com/lsernau/grh-lean/blob/72fead102f663d5c7a407964a13dc0cf9d977207/grh-compact/src/New/Seesaw/PairTensorData.lean\#L161-L163}{{\leanfont Data.\allowbreak{}descentMove}}}
\expandafter\newcommand\csname LN@liftMove\endcsname{\href{https://github.com/lsernau/grh-lean/blob/72fead102f663d5c7a407964a13dc0cf9d977207/grh-compact/src/New/Seesaw/PairTensorData.lean\#L167-L169}{{\leanfont Data.\allowbreak{}liftMove}}}
\expandafter\newcommand\csname LN@gridChardet\endcsname{\href{https://github.com/lsernau/grh-lean/blob/72fead102f663d5c7a407964a13dc0cf9d977207/grh-compact/src/New/Seesaw/PairTensorData.lean\#L173-L177}{{\leanfont Data.\allowbreak{}gridChar\_det}}}
\expandafter\newcommand\csname LN@levelIndexequntwisted\endcsname{\href{https://github.com/lsernau/grh-lean/blob/72fead102f663d5c7a407964a13dc0cf9d977207/grh-compact/src/New/Seesaw/PairTensorData.lean\#L181-L184}{{\leanfont Data.\allowbreak{}levelIndex\_eq\_untwisted}}}
\newcommand{\ConcordanceRows}{\crow{hasLogBoundvarpi}{\href{https://github.com/lsernau/grh-lean/blob/72fead102f663d5c7a407964a13dc0cf9d977207/grh-compact/src/New/Seesaw/IndexTypes.lean\#L19-L27}{{\leanfont Data.\allowbreak{}hasLogBound\_varpi}}}{IndexTypes.lean}{19--27}
\crow{hasLogBoundrootDisc}{\href{https://github.com/lsernau/grh-lean/blob/72fead102f663d5c7a407964a13dc0cf9d977207/grh-compact/src/New/Seesaw/IndexTypes.lean\#L30-L33}{{\leanfont Data.\allowbreak{}hasLogBound\_rootDisc}}}{IndexTypes.lean}{30--33}
\crow{piGBar}{\href{https://github.com/lsernau/grh-lean/blob/72fead102f663d5c7a407964a13dc0cf9d977207/grh-compact/src/New/Seesaw/IndexTypes.lean\#L37-L37}{{\leanfont Data.\allowbreak{}piGBar}}}{IndexTypes.lean}{37--37}
\crow{IsChosen}{\href{https://github.com/lsernau/grh-lean/blob/72fead102f663d5c7a407964a13dc0cf9d977207/grh-compact/src/New/Seesaw/IndexTypes.lean\#L41-L41}{{\leanfont Data.\allowbreak{}IsChosen}}}{IndexTypes.lean}{41--41}
\crow{ChosenPlace}{\href{https://github.com/lsernau/grh-lean/blob/72fead102f663d5c7a407964a13dc0cf9d977207/grh-compact/src/New/Seesaw/IndexTypes.lean\#L45-L45}{{\leanfont Data.\allowbreak{}ChosenPlace}}}{IndexTypes.lean}{45--45}
\crow{piGBarmulstar}{\href{https://github.com/lsernau/grh-lean/blob/72fead102f663d5c7a407964a13dc0cf9d977207/grh-compact/src/New/Seesaw/IndexTypes.lean\#L47-L48}{{\leanfont Data.\allowbreak{}piGBar\_mul\_star}}}{IndexTypes.lean}{47--48}
\crow{piGnezero}{\href{https://github.com/lsernau/grh-lean/blob/72fead102f663d5c7a407964a13dc0cf9d977207/grh-compact/src/New/Seesaw/IndexTypes.lean\#L50-L50}{{\leanfont Data.\allowbreak{}piG\_ne\_zero}}}{IndexTypes.lean}{50--50}
\crow{piGBarnezero}{\href{https://github.com/lsernau/grh-lean/blob/72fead102f663d5c7a407964a13dc0cf9d977207/grh-compact/src/New/Seesaw/IndexTypes.lean\#L52-L52}{{\leanfont Data.\allowbreak{}piGBar\_ne\_zero}}}{IndexTypes.lean}{52--52}
\crow{existscofactor}{\href{https://github.com/lsernau/grh-lean/blob/72fead102f663d5c7a407964a13dc0cf9d977207/grh-compact/src/New/Seesaw/IndexTypes.lean\#L54-L55}{{\leanfont Data.\allowbreak{}exists\_cofactor}}}{IndexTypes.lean}{54--55}
\crow{isChoseniff}{\href{https://github.com/lsernau/grh-lean/blob/72fead102f663d5c7a407964a13dc0cf9d977207/grh-compact/src/New/Seesaw/IndexTypes.lean\#L57-L58}{{\leanfont Data.\allowbreak{}isChosen\_iff}}}{IndexTypes.lean}{57--58}
\crow{qeqg}{\href{https://github.com/lsernau/grh-lean/blob/72fead102f663d5c7a407964a13dc0cf9d977207/grh-compact/src/New/Seesaw/IndexTypes.lean\#L62-L63}{{\leanfont Data.\allowbreak{}q\_eq\_g}}}{IndexTypes.lean}{62--63}
\crow{finiteisChosen}{\href{https://github.com/lsernau/grh-lean/blob/72fead102f663d5c7a407964a13dc0cf9d977207/grh-compact/src/New/Seesaw/IndexTypes.lean\#L65-L66}{{\leanfont Data.\allowbreak{}finite\_isChosen}}}{IndexTypes.lean}{65--66}
\crow{existsresiduecharoforder}{\href{https://github.com/lsernau/grh-lean/blob/72fead102f663d5c7a407964a13dc0cf9d977207/grh-compact/src/New/Seesaw/IndexTypes.lean\#L70-L72}{{\leanfont Data.\allowbreak{}exists\_residue\_char\_of\_order}}}{IndexTypes.lean}{70--72}
\crow{decla}{\href{https://github.com/lsernau/grh-lean/blob/72fead102f663d5c7a407964a13dc0cf9d977207/grh-compact/src/New/Seesaw/IndexTypes.lean\#L76-L76}{{\leanfont Data.\allowbreak{}ω}}}{IndexTypes.lean}{76--76}
\crow{orderOf}{\href{https://github.com/lsernau/grh-lean/blob/72fead102f663d5c7a407964a13dc0cf9d977207/grh-compact/src/New/Seesaw/IndexTypes.lean\#L79-L80}{{\leanfont Data.\allowbreak{}orderOf\_ω}}}{IndexTypes.lean}{79--80}
\crow{omegaeqoneofnotisChosen}{\href{https://github.com/lsernau/grh-lean/blob/72fead102f663d5c7a407964a13dc0cf9d977207/grh-compact/src/New/Seesaw/IndexTypes.lean\#L82-L82}{{\leanfont Data.\allowbreak{}omega\_eq\_one\_of\_not\_isChosen}}}{IndexTypes.lean}{82--82}
\crow{omegazpowmulpoweqone}{\href{https://github.com/lsernau/grh-lean/blob/72fead102f663d5c7a407964a13dc0cf9d977207/grh-compact/src/New/Seesaw/IndexTypes.lean\#L84-L85}{{\leanfont Data.\allowbreak{}omega\_zpow\_mul\_pow\_eq\_one}}}{IndexTypes.lean}{84--85}
\crow{existscharhasWeight}{\href{https://github.com/lsernau/grh-lean/blob/72fead102f663d5c7a407964a13dc0cf9d977207/grh-compact/src/New/Seesaw/IndexTypes.lean\#L89-L97}{{\leanfont Data.\allowbreak{}exists\_char\_hasWeight}}}{IndexTypes.lean}{89--97}
\crow{muFlat}{\href{https://github.com/lsernau/grh-lean/blob/72fead102f663d5c7a407964a13dc0cf9d977207/grh-compact/src/New/Seesaw/IndexTypes.lean\#L101-L101}{{\leanfont Data.\allowbreak{}muFlat}}}{IndexTypes.lean}{101--101}
\crow{declb}{\href{https://github.com/lsernau/grh-lean/blob/72fead102f663d5c7a407964a13dc0cf9d977207/grh-compact/src/New/Seesaw/IndexTypes.lean\#L105-L105}{{\leanfont Data.\allowbreak{}χ}}}{IndexTypes.lean}{105--105}
\crow{muFlathasWeight}{\href{https://github.com/lsernau/grh-lean/blob/72fead102f663d5c7a407964a13dc0cf9d977207/grh-compact/src/New/Seesaw/IndexTypes.lean\#L108-L108}{{\leanfont Data.\allowbreak{}muFlat\_hasWeight}}}{IndexTypes.lean}{108--108}
\crow{chihasWeight}{\href{https://github.com/lsernau/grh-lean/blob/72fead102f663d5c7a407964a13dc0cf9d977207/grh-compact/src/New/Seesaw/IndexTypes.lean\#L111-L111}{{\leanfont Data.\allowbreak{}chi\_hasWeight}}}{IndexTypes.lean}{111--111}
\crow{Pair}{\href{https://github.com/lsernau/grh-lean/blob/72fead102f663d5c7a407964a13dc0cf9d977207/grh-compact/src/New/Seesaw/IndexTypes.lean\#L115-L122}{{\leanfont Data.\allowbreak{}Pair}}}{IndexTypes.lean}{115--122}
\crow{rowParameter}{\href{https://github.com/lsernau/grh-lean/blob/72fead102f663d5c7a407964a13dc0cf9d977207/grh-compact/src/New/Seesaw/IndexTypes.lean\#L126-L126}{{\leanfont Data.\allowbreak{}rowParameter}}}{IndexTypes.lean}{126--126}
\crow{logLevelTarget}{\href{https://github.com/lsernau/grh-lean/blob/72fead102f663d5c7a407964a13dc0cf9d977207/grh-compact/src/New/Seesaw/IndexTypes.lean\#L133-L134}{{\leanfont Data.\allowbreak{}logLevelTarget}}}{IndexTypes.lean}{133--134}
\crow{logLevelTargetmemIcc}{\href{https://github.com/lsernau/grh-lean/blob/72fead102f663d5c7a407964a13dc0cf9d977207/grh-compact/src/New/Seesaw/IndexTypes.lean\#L137-L141}{{\leanfont Data.\allowbreak{}logLevelTarget\_mem\_Icc}}}{IndexTypes.lean}{137--141}
\crow{N}{\href{https://github.com/lsernau/grh-lean/blob/72fead102f663d5c7a407964a13dc0cf9d977207/grh-compact/src/New/Seesaw/IndexTypes.lean\#L145-L146}{{\leanfont Data.\allowbreak{}N}}}{IndexTypes.lean}{145--146}
\crow{Nprime}{\href{https://github.com/lsernau/grh-lean/blob/72fead102f663d5c7a407964a13dc0cf9d977207/grh-compact/src/New/Seesaw/IndexTypes.lean\#L149-L149}{{\leanfont Data.\allowbreak{}N\_prime}}}{IndexTypes.lean}{149--149}
\crow{NmodEqone}{\href{https://github.com/lsernau/grh-lean/blob/72fead102f663d5c7a407964a13dc0cf9d977207/grh-compact/src/New/Seesaw/IndexTypes.lean\#L151-L152}{{\leanfont Data.\allowbreak{}N\_modEq\_one}}}{IndexTypes.lean}{151--152}
\crow{logNmemwindow}{\href{https://github.com/lsernau/grh-lean/blob/72fead102f663d5c7a407964a13dc0cf9d977207/grh-compact/src/New/Seesaw/IndexTypes.lean\#L157-L158}{{\leanfont Data.\allowbreak{}log\_N\_mem\_window}}}{IndexTypes.lean}{157--158}
\crow{logNmemIcc}{\href{https://github.com/lsernau/grh-lean/blob/72fead102f663d5c7a407964a13dc0cf9d977207/grh-compact/src/New/Seesaw/IndexTypes.lean\#L161-L162}{{\leanfont Data.\allowbreak{}log\_N\_mem\_Icc}}}{IndexTypes.lean}{161--162}
\crow{existspimulstar}{\href{https://github.com/lsernau/grh-lean/blob/72fead102f663d5c7a407964a13dc0cf9d977207/grh-compact/src/New/Seesaw/IndexTypes.lean\#L173-L175}{{\leanfont Data.\allowbreak{}exists\_pi\_mul\_star}}}{IndexTypes.lean}{173--175}
\crow{declc}{\href{https://github.com/lsernau/grh-lean/blob/72fead102f663d5c7a407964a13dc0cf9d977207/grh-compact/src/New/Seesaw/IndexTypes.lean\#L178-L178}{{\leanfont Data.\allowbreak{}π}}}{IndexTypes.lean}{178--178}
\crow{mulstar}{\href{https://github.com/lsernau/grh-lean/blob/72fead102f663d5c7a407964a13dc0cf9d977207/grh-compact/src/New/Seesaw/IndexTypes.lean\#L182-L182}{{\leanfont Data.\allowbreak{}π\_mul\_star}}}{IndexTypes.lean}{182--182}
\crow{TameChar}{\href{https://github.com/lsernau/grh-lean/blob/72fead102f663d5c7a407964a13dc0cf9d977207/grh-compact/src/New/Seesaw/IndexTypes.lean\#L186-L192}{{\leanfont Data.\allowbreak{}TameChar}}}{IndexTypes.lean}{186--192}
\crow{levelNorm}{\href{https://github.com/lsernau/grh-lean/blob/72fead102f663d5c7a407964a13dc0cf9d977207/grh-compact/src/New/Seesaw/IndexTypes.lean\#L195-L195}{{\leanfont Data.\allowbreak{}TameChar.\allowbreak{}levelNorm}}}{IndexTypes.lean}{195--195}
\crow{Twist}{\href{https://github.com/lsernau/grh-lean/blob/72fead102f663d5c7a407964a13dc0cf9d977207/grh-compact/src/New/Seesaw/IndexTypes.lean\#L199-L199}{{\leanfont Data.\allowbreak{}Twist}}}{IndexTypes.lean}{199--199}
\crow{one}{\href{https://github.com/lsernau/grh-lean/blob/72fead102f663d5c7a407964a13dc0cf9d977207/grh-compact/src/New/Seesaw/IndexTypes.lean\#L206-L206}{{\leanfont Data.\allowbreak{}TameChar.\allowbreak{}one}}}{IndexTypes.lean}{206--206}
\crow{decld}{\href{https://github.com/lsernau/grh-lean/blob/72fead102f663d5c7a407964a13dc0cf9d977207/grh-compact/src/New/Seesaw/IndexTypes.lean\#L210-L210}{{\leanfont Data.\allowbreak{}μ₀}}}{IndexTypes.lean}{210--210}
\crow{muhasWeight}{\href{https://github.com/lsernau/grh-lean/blob/72fead102f663d5c7a407964a13dc0cf9d977207/grh-compact/src/New/Seesaw/IndexTypes.lean\#L213-L215}{{\leanfont Data.\allowbreak{}mu₀\_hasWeight}}}{IndexTypes.lean}{213--215}
\crow{liftChar}{\href{https://github.com/lsernau/grh-lean/blob/72fead102f663d5c7a407964a13dc0cf9d977207/grh-compact/src/New/Seesaw/IndexTypes.lean\#L218-L218}{{\leanfont Data.\allowbreak{}liftChar}}}{IndexTypes.lean}{218--218}
\crow{liftTamePart}{\href{https://github.com/lsernau/grh-lean/blob/72fead102f663d5c7a407964a13dc0cf9d977207/grh-compact/src/New/Seesaw/IndexTypes.lean\#L222-L222}{{\leanfont Data.\allowbreak{}liftTamePart}}}{IndexTypes.lean}{222--222}
\crow{hasWeightliftChar}{\href{https://github.com/lsernau/grh-lean/blob/72fead102f663d5c7a407964a13dc0cf9d977207/grh-compact/src/New/Seesaw/IndexTypes.lean\#L224-L225}{{\leanfont Data.\allowbreak{}hasWeight\_liftChar}}}{IndexTypes.lean}{224--225}
\crow{LiftIndex}{\href{https://github.com/lsernau/grh-lean/blob/72fead102f663d5c7a407964a13dc0cf9d977207/grh-compact/src/New/Seesaw/IndexTypes.lean\#L228-L228}{{\leanfont Data.\allowbreak{}LiftIndex}}}{IndexTypes.lean}{228--228}
\crow{BlockIndex}{\href{https://github.com/lsernau/grh-lean/blob/72fead102f663d5c7a407964a13dc0cf9d977207/grh-compact/src/New/Seesaw/IndexTypes.lean\#L231-L231}{{\leanfont BlockIndex}}}{IndexTypes.lean}{231--231}
\crow{firstBlock}{\href{https://github.com/lsernau/grh-lean/blob/72fead102f663d5c7a407964a13dc0cf9d977207/grh-compact/src/New/Seesaw/IndexTypes.lean\#L234-L235}{{\leanfont BlockIndex.\allowbreak{}firstBlock}}}{IndexTypes.lean}{234--235}
\crow{equivSum}{\href{https://github.com/lsernau/grh-lean/blob/72fead102f663d5c7a407964a13dc0cf9d977207/grh-compact/src/New/Seesaw/IndexTypes.lean\#L239-L241}{{\leanfont BlockIndex.\allowbreak{}equivSum}}}{IndexTypes.lean}{239--241}
\crow{instFintypeYT}{\href{https://github.com/lsernau/grh-lean/blob/72fead102f663d5c7a407964a13dc0cf9d977207/grh-compact/src/New/Seesaw/IndexTypes.lean\#L245-L246}{{\leanfont instFintypeYT}}}{IndexTypes.lean}{245--246}
\crow{card}{\href{https://github.com/lsernau/grh-lean/blob/72fead102f663d5c7a407964a13dc0cf9d977207/grh-compact/src/New/Seesaw/IndexTypes.lean\#L248-L249}{{\leanfont BlockIndex.\allowbreak{}card}}}{IndexTypes.lean}{248--249}
\crow{decle}{\href{https://github.com/lsernau/grh-lean/blob/72fead102f663d5c7a407964a13dc0cf9d977207/grh-compact/src/New/Seesaw/IndexTypes.lean\#L252-L252}{{\leanfont Data.\allowbreak{}ϱ}}}{IndexTypes.lean}{252--252}
\crow{InitIndex}{\href{https://github.com/lsernau/grh-lean/blob/72fead102f663d5c7a407964a13dc0cf9d977207/grh-compact/src/New/Seesaw/IndexTypes.lean\#L255-L255}{{\leanfont Data.\allowbreak{}InitIndex}}}{IndexTypes.lean}{255--255}
\crow{X}{\href{https://github.com/lsernau/grh-lean/blob/72fead102f663d5c7a407964a13dc0cf9d977207/grh-compact/src/New/Seesaw/IndexTypes.lean\#L259-L259}{{\leanfont Data.\allowbreak{}X}}}{IndexTypes.lean}{259--259}
\crow{shape}{\href{https://github.com/lsernau/grh-lean/blob/72fead102f663d5c7a407964a13dc0cf9d977207/grh-compact/src/New/Seesaw/IndexTypes.lean\#L263-L263}{{\leanfont Data.\allowbreak{}shape}}}{IndexTypes.lean}{263--263}
\crow{RowIndex}{\href{https://github.com/lsernau/grh-lean/blob/72fead102f663d5c7a407964a13dc0cf9d977207/grh-compact/src/New/Seesaw/IndexTypes.lean\#L267-L267}{{\leanfont Data.\allowbreak{}RowIndex}}}{IndexTypes.lean}{267--267}
\crow{Y}{\href{https://github.com/lsernau/grh-lean/blob/72fead102f663d5c7a407964a13dc0cf9d977207/grh-compact/src/New/Seesaw/IndexTypes.lean\#L270-L270}{{\leanfont Data.\allowbreak{}Y}}}{IndexTypes.lean}{270--270}
\crow{cardrowIndex}{\href{https://github.com/lsernau/grh-lean/blob/72fead102f663d5c7a407964a13dc0cf9d977207/grh-compact/src/New/Seesaw/IndexTypes.lean\#L274-L274}{{\leanfont Data.\allowbreak{}card\_rowIndex}}}{IndexTypes.lean}{274--274}
\crow{basisNumber}{\href{https://github.com/lsernau/grh-lean/blob/72fead102f663d5c7a407964a13dc0cf9d977207/grh-compact/src/New/Seesaw/IndexTypes.lean\#L281-L281}{{\leanfont Data.\allowbreak{}basisNumber}}}{IndexTypes.lean}{281--281}
\crow{paramW}{\href{https://github.com/lsernau/grh-lean/blob/72fead102f663d5c7a407964a13dc0cf9d977207/grh-compact/src/New/Seesaw/IndexTypes.lean\#L285-L285}{{\leanfont Data.\allowbreak{}paramW}}}{IndexTypes.lean}{285--285}
\crow{entrySign}{\href{https://github.com/lsernau/grh-lean/blob/72fead102f663d5c7a407964a13dc0cf9d977207/grh-compact/src/New/Seesaw/IndexTypes.lean\#L289-L289}{{\leanfont Data.\allowbreak{}entrySign}}}{IndexTypes.lean}{289--289}
\crow{iteoddeqsubemod}{\href{https://github.com/lsernau/grh-lean/blob/72fead102f663d5c7a407964a13dc0cf9d977207/grh-compact/src/New/Seesaw/IndexTypes.lean\#L291-L291}{{\leanfont Data.\allowbreak{}ite\_odd\_eq\_sub\_emod}}}{IndexTypes.lean}{291--291}
\crow{signsW}{\href{https://github.com/lsernau/grh-lean/blob/72fead102f663d5c7a407964a13dc0cf9d977207/grh-compact/src/New/Seesaw/IndexTypes.lean\#L295-L295}{{\leanfont Data.\allowbreak{}signsW}}}{IndexTypes.lean}{295--295}
\crow{chosen}{\href{https://github.com/lsernau/grh-lean/blob/72fead102f663d5c7a407964a13dc0cf9d977207/grh-compact/src/New/Seesaw/IndexTypes.lean\#L299-L299}{{\leanfont Data.\allowbreak{}chosen}}}{IndexTypes.lean}{299--299}
\crow{G}{\href{https://github.com/lsernau/grh-lean/blob/72fead102f663d5c7a407964a13dc0cf9d977207/grh-compact/src/New/Seesaw/IndexTypes.lean\#L302-L302}{{\leanfont Data.\allowbreak{}G}}}{IndexTypes.lean}{302--302}
\crow{goodTupleFamilies}{\href{https://github.com/lsernau/grh-lean/blob/72fead102f663d5c7a407964a13dc0cf9d977207/grh-compact/src/New/Seesaw/IndexTypes.lean\#L308-L310}{{\leanfont Data.\allowbreak{}goodTupleFamilies}}}{IndexTypes.lean}{308--310}
\crow{e}{\href{https://github.com/lsernau/grh-lean/blob/72fead102f663d5c7a407964a13dc0cf9d977207/grh-compact/src/New/Seesaw/IndexTypes.lean\#L314-L314}{{\leanfont Data.\allowbreak{}e₀}}}{IndexTypes.lean}{314--314}
\crow{existsbranchingpairing}{\href{https://github.com/lsernau/grh-lean/blob/72fead102f663d5c7a407964a13dc0cf9d977207/grh-compact/src/New/Seesaw/PairTensorData.lean\#L22-L58}{{\leanfont exists\_branching\_pairing}}}{PairTensorData.lean}{22--58}
\crow{gridSpace}{\href{https://github.com/lsernau/grh-lean/blob/72fead102f663d5c7a407964a13dc0cf9d977207/grh-compact/src/New/Seesaw/PairTensorData.lean\#L68-L71}{{\leanfont Data.\allowbreak{}gridSpace}}}{PairTensorData.lean}{68--71}
\crow{gridKernel}{\href{https://github.com/lsernau/grh-lean/blob/72fead102f663d5c7a407964a13dc0cf9d977207/grh-compact/src/New/Seesaw/PairTensorData.lean\#L75-L75}{{\leanfont Data.\allowbreak{}gridKernel}}}{PairTensorData.lean}{75--75}
\crow{GridForm}{\href{https://github.com/lsernau/grh-lean/blob/72fead102f663d5c7a407964a13dc0cf9d977207/grh-compact/src/New/Seesaw/PairTensorData.lean\#L79-L79}{{\leanfont Data.\allowbreak{}GridForm}}}{PairTensorData.lean}{79--79}
\crow{gridChar}{\href{https://github.com/lsernau/grh-lean/blob/72fead102f663d5c7a407964a13dc0cf9d977207/grh-compact/src/New/Seesaw/PairTensorData.lean\#L83-L84}{{\leanfont Data.\allowbreak{}gridChar}}}{PairTensorData.lean}{83--84}
\crow{gridTamePart}{\href{https://github.com/lsernau/grh-lean/blob/72fead102f663d5c7a407964a13dc0cf9d977207/grh-compact/src/New/Seesaw/PairTensorData.lean\#L88-L88}{{\leanfont Data.\allowbreak{}gridTamePart}}}{PairTensorData.lean}{88--88}
\crow{triangleOrder}{\href{https://github.com/lsernau/grh-lean/blob/72fead102f663d5c7a407964a13dc0cf9d977207/grh-compact/src/New/Seesaw/PairTensorData.lean\#L92-L92}{{\leanfont Data.\allowbreak{}triangleOrder}}}{PairTensorData.lean}{92--92}
\crow{Node}{\href{https://github.com/lsernau/grh-lean/blob/72fead102f663d5c7a407964a13dc0cf9d977207/grh-compact/src/New/Seesaw/PairTensorData.lean\#L96-L96}{{\leanfont Data.\allowbreak{}Node}}}{PairTensorData.lean}{96--96}
\crow{NodeIndex}{\href{https://github.com/lsernau/grh-lean/blob/72fead102f663d5c7a407964a13dc0cf9d977207/grh-compact/src/New/Seesaw/PairTensorData.lean\#L99-L99}{{\leanfont Data.\allowbreak{}NodeIndex}}}{PairTensorData.lean}{99--99}
\crow{nodeLevel}{\href{https://github.com/lsernau/grh-lean/blob/72fead102f663d5c7a407964a13dc0cf9d977207/grh-compact/src/New/Seesaw/PairTensorData.lean\#L103-L104}{{\leanfont Data.\allowbreak{}nodeLevel}}}{PairTensorData.lean}{103--104}
\crow{levelIndex}{\href{https://github.com/lsernau/grh-lean/blob/72fead102f663d5c7a407964a13dc0cf9d977207/grh-compact/src/New/Seesaw/PairTensorData.lean\#L108-L108}{{\leanfont Data.\allowbreak{}levelIndex}}}{PairTensorData.lean}{108--108}
\crow{gridLevelForms}{\href{https://github.com/lsernau/grh-lean/blob/72fead102f663d5c7a407964a13dc0cf9d977207/grh-compact/src/New/Seesaw/PairTensorData.lean\#L112-L113}{{\leanfont Data.\allowbreak{}gridLevelForms}}}{PairTensorData.lean}{112--113}
\crow{isConjKernelPairgridKernel}{\href{https://github.com/lsernau/grh-lean/blob/72fead102f663d5c7a407964a13dc0cf9d977207/grh-compact/src/New/Seesaw/PairTensorData.lean\#L115-L117}{{\leanfont Data.\allowbreak{}isConjKernelPair\_gridKernel}}}{PairTensorData.lean}{115--117}
\crow{gridCharhasWeight}{\href{https://github.com/lsernau/grh-lean/blob/72fead102f663d5c7a407964a13dc0cf9d977207/grh-compact/src/New/Seesaw/PairTensorData.lean\#L119-L120}{{\leanfont Data.\allowbreak{}gridChar\_hasWeight}}}{PairTensorData.lean}{119--120}
\crow{isEdgeFamilygrid}{\href{https://github.com/lsernau/grh-lean/blob/72fead102f663d5c7a407964a13dc0cf9d977207/grh-compact/src/New/Seesaw/PairTensorData.lean\#L125-L128}{{\leanfont Data.\allowbreak{}isEdgeFamily\_grid}}}{PairTensorData.lean}{125--128}
\crow{gridSpacediag}{\href{https://github.com/lsernau/grh-lean/blob/72fead102f663d5c7a407964a13dc0cf9d977207/grh-compact/src/New/Seesaw/PairTensorData.lean\#L132-L136}{{\leanfont Data.\allowbreak{}gridSpace\_diag}}}{PairTensorData.lean}{132--136}
\crow{gridSpacesubdiag}{\href{https://github.com/lsernau/grh-lean/blob/72fead102f663d5c7a407964a13dc0cf9d977207/grh-compact/src/New/Seesaw/PairTensorData.lean\#L140-L143}{{\leanfont Data.\allowbreak{}gridSpace\_subdiag}}}{PairTensorData.lean}{140--143}
\crow{pointDescent}{\href{https://github.com/lsernau/grh-lean/blob/72fead102f663d5c7a407964a13dc0cf9d977207/grh-compact/src/New/Seesaw/PairTensorData.lean\#L148-L150}{{\leanfont Data.\allowbreak{}pointDescent}}}{PairTensorData.lean}{148--150}
\crow{iotaPlus}{\href{https://github.com/lsernau/grh-lean/blob/72fead102f663d5c7a407964a13dc0cf9d977207/grh-compact/src/New/Seesaw/PairTensorData.lean\#L154-L155}{{\leanfont Data.\allowbreak{}iotaPlus}}}{PairTensorData.lean}{154--155}
\crow{descentMove}{\href{https://github.com/lsernau/grh-lean/blob/72fead102f663d5c7a407964a13dc0cf9d977207/grh-compact/src/New/Seesaw/PairTensorData.lean\#L161-L163}{{\leanfont Data.\allowbreak{}descentMove}}}{PairTensorData.lean}{161--163}
\crow{liftMove}{\href{https://github.com/lsernau/grh-lean/blob/72fead102f663d5c7a407964a13dc0cf9d977207/grh-compact/src/New/Seesaw/PairTensorData.lean\#L167-L169}{{\leanfont Data.\allowbreak{}liftMove}}}{PairTensorData.lean}{167--169}
\crow{gridChardet}{\href{https://github.com/lsernau/grh-lean/blob/72fead102f663d5c7a407964a13dc0cf9d977207/grh-compact/src/New/Seesaw/PairTensorData.lean\#L173-L177}{{\leanfont Data.\allowbreak{}gridChar\_det}}}{PairTensorData.lean}{173--177}
\crow{levelIndexequntwisted}{\href{https://github.com/lsernau/grh-lean/blob/72fead102f663d5c7a407964a13dc0cf9d977207/grh-compact/src/New/Seesaw/PairTensorData.lean\#L181-L184}{{\leanfont Data.\allowbreak{}levelIndex\_eq\_untwisted}}}{PairTensorData.lean}{181--184}}
% the notations of the defined functions: one slot for each explicit argument of the Lean declaration
\newcommand{\piGbar}[1]{\bar\pi_g(#1)}
\newcommand{\isChosen}[2]{#2\in\mathrm{Ch}(#1)}
\newcommand{\notChosen}[2]{#2\notin\mathrm{Ch}(#1)}
\newcommand{\ChPl}[2]{\mathrm{Ch}_{#2}(#1)}
\newcommand{\omD}[2]{\omega_{#1,#2}}
\newcommand{\muFlat}[1]{\mu^{\flat}_{#1}}
\newcommand{\chiCirc}[1]{\chi^{\circ}_{#1}}
\newcommand{\PairSet}[1]{\mathrm{Pair}(#1)}
\newcommand{\rowPar}[3]{\mathrm{par}_{#1,#2}(#3)}
\newcommand{\logTarget}[1]{\mathrm{tar}(#1)}
\newcommand{\levelN}[1]{N(#1)}
\newcommand{\piN}[1]{\pi_N(#1)}
\newcommand{\TameSet}[1]{\mathrm{Tame}(#1)}
\newcommand{\levelNorm}[1]{\mathrm{n}(#1)}
\newcommand{\TwistSet}[1]{\mathrm{Tw}(#1)}
\newcommand{\muZero}[2]{\mu^{0}_{#1,#2}}
\newcommand{\liftChar}[3]{\mu_{#1,#3}(#2)}
\newcommand{\liftTame}[3]{\psi_{#1,#3}(#2)}
\newcommand{\LiftIdx}[1]{\mathrm{Lf}(#1)}
\newcommand{\BlockIdx}[3]{\operatorname{Bl}_{#1}(#2,#3)}
\newcommand{\firstBlock}[1]{\operatorname{in}_{#1}}
\newcommand{\equivSum}[2]{\operatorname{eq}_{#2}(#1)}
\newcommand{\rhoD}[1]{\mathrm{hm}(#1)}
\newcommand{\Xidx}[2]{\underline{X}_{#2}(#1)}
\newcommand{\shapeD}[1]{\mathfrak{s}(#1)}
\newcommand{\Ridx}[1]{\underline{R}(#1)}
\newcommand{\Yidx}[2]{\underline{Y}_{#2}(#1)}
\newcommand{\basisNum}[1]{\operatorname{num}_{#1}}
\newcommand{\parW}[3]{\mathrm{par}^{W}_{#1,#2}(#3)}
\newcommand{\signsW}[3]{\operatorname{sg}^{W}_{#1,#2}(#3)}
\newcommand{\chosenOf}[2]{\operatorname{pl}_{#1}(#2)}
\newcommand{\Ggrp}[1]{\mathsf{G}(#1)}
\newcommand{\goodFam}[1]{\mathrm{Good}(#1)}
\newcommand{\ezero}[1]{e_0(#1)}
\newcommand{\gridSp}[4]{\mathcal{V}^{#2}_{#3,#4}[#1](D)}
\newcommand{\gridKer}[3]{\mathrm{ker}^{#1}_{#2,#3}(D)}
\newcommand{\gridForm}[3]{\mathcal{F}^{#1}_{#2,#3}(D)}
\newcommand{\gridChar}[4]{\mu^{\mathrm{gr}}_{D,#3}(#1;#2,#4)}
\newcommand{\gridTame}[3]{\psi^{\mathrm{gr}}_{D,#3}(#1;#2)}
\newcommand{\triOrd}[3]{\operatorname{to}_{D}(#1,#2,#3)}
\newcommand{\nodeIdx}[1]{\underline{\mathcal{V}}_{D}(#1)}
\newcommand{\levIdx}[4]{\operatorname{ix}_{D}(#1;#2,#3,#4)}
\newcommand{\gridLevForms}[1]{\mathcal{F}^{\mathrm{lev}}_{#1}(D)}
\newcommand{\ptDesc}[4]{\operatorname{desc}^{#2}_{D,#3}(#1)(#4)}
\newcommand{\iotaPlus}[4]{\iota^{+,#2}_{D,#3}(#1)(#4)}
\newcommand{\descMove}[2]{\operatorname{desc}_{D}[#1](#2)}
\newcommand{\liftMove}[2]{\operatorname{lift}_{D}[#1](#2)}
\newcommand{\UnitMembers}{312}
\newcommand{\UnitInNote}{88}
\newcommand{\UnitProofs}{83}
\newcommand{\UnitSteps}{503}
\newcommand{\UnitAllSteps}{550}


\linespread{1.08}                       % 11 pt on 14.7 pt
\setlength{\parindent}{0pt}
\setlength{\parskip}{5.5pt plus 1.5pt minus 0.5pt}
\setlength{\abovedisplayskip}{7.5pt plus 2pt minus 3pt}
\setlength{\belowdisplayskip}{7.5pt plus 2pt minus 3pt}
\setlength{\abovedisplayshortskip}{3pt plus 2pt}
\setlength{\belowdisplayshortskip}{6pt plus 2pt minus 2pt}
\emergencystretch=1.5em
\clubpenalty=10000 \widowpenalty=10000 \displaywidowpenalty=10000
\relpenalty=7000 \binoppenalty=9000
\raggedbottom
\pagestyle{fancy}\fancyhf{}\renewcommand{\headrulewidth}{0pt}
\fancyfoot[C]{\sffamily\fontsize{7.6}{9}\selectfont\color{inkfaint}\thepage\,/\,\pageref*{LastPage}}
\setlist{topsep=2pt,partopsep=0pt,parsep=0pt,itemsep=4pt,leftmargin=2.1em,labelsep=0.55em}

\newcommand{\Q}{\mathbb{Q}}
\newcommand{\Z}{\mathbb{Z}}
\newcommand{\F}{\mathbb{F}}
\newcommand{\len}{\operatorname{len}}
\newcommand{\cok}{\operatorname{cok}}
\newcommand{\End}{\operatorname{End}}
\newcommand{\Hom}{\operatorname{Hom}}
\newcommand{\Gal}{\operatorname{Gal}}
\newcommand{\diag}{\operatorname{diag}}
\newcommand{\kw}[1]{\emph{#1}}
\newcommand{\m}[1]{\mbox{$#1$}}        % a short formula that must not be broken across lines
\newcommand{\head}[1]{\par\addvspace{17pt plus 4pt minus 2pt}\Needspace{6\baselineskip}%
  \noindent{\fontsize{13.2}{16}\selectfont\bfseries\color{ink}#1}\par\nopagebreak\vspace{-1pt}}
\newcommand{\sect}[2]{\refstepcounter{section}\label{#2}\par\addvspace{19pt plus 4pt minus 2pt}\Needspace{7\baselineskip}%
  \noindent\llap{{\fontsize{13.2}{16}\selectfont\bfseries\color{inkfaint}\thesection}\hspace{9pt}}%
  {\fontsize{13.2}{16}\selectfont\bfseries\color{ink}#1}\par\nopagebreak\vspace{-1pt}}
\newcommand{\subhead}[1]{\par\addvspace{9pt plus 2pt}\Needspace{4\baselineskip}\noindent\textbf{#1}\enspace\ignorespaces}
% where a statement lives in the Lean archive
\newcommand{\leantag}{{\sffamily\bfseries\fontsize{6.3}{8}\selectfont\color{leancol}\addfontfeatures{LetterSpace=7}LEAN}}
\newcommand{\sep}{\unskip\nobreak\hspace{5pt}{\color{inkfaint}\textperiodcentered}\hspace{5pt plus 2pt}\ignorespaces}
\newcommand{\leanline}[1]{\par\nopagebreak\vspace{-2.5pt}{\fontsize{8.9}{12.4}\selectfont\color{inksoft}\raggedright
  \noindent\llap{\leantag\hspace{9pt}}#1\par}\vspace{1pt}}
% a displayed statement in a tinted box: #1 colour, #2 text
\newcommand{\statement}[2]{\par\addvspace{7pt}\begingroup\setlength{\fboxsep}{8.5pt}\setlength{\fboxrule}{0.6pt}%
  \noindent\fcolorbox{#1!30!white}{#1!3!white}{\parbox{\dimexpr\linewidth-2\fboxsep-2\fboxrule\relax}{#2}}%
  \endgroup\par\addvspace{7pt}}
\newcommand{\nt}[2]{\par\noindent\hangindent=1.1em\hangafter=1 #1\enspace{\color{inksoft}#2}\par\vspace{1pt}}
\newcommand{\legdot}[1]{\tikz[baseline=-0.62ex]\fill[#1] (0,0) circle (2.5pt);\hspace{3pt}}
\newcommand{\crow}[4]{\csname CS@#1\endcsname & #2 & \textcolor{inksoft}{#3} & \textcolor{inksoft}{#4}\\}
\newcommand{\leanquoteline}[2]{\par\noindent\hangindent=\dimexpr#1\fontdimen2\font+2.2em\relax\hangafter=1
  \hspace*{#1\fontdimen2\font}#2}
\newcommand{\quoted}[1]{\par\vspace{3pt}{\ttfamily\fontsize{7.5}{10.4}\selectfont\color{leancol}\raggedright
  \setlength{\parskip}{0pt}\leanquote{#1}\par}\vspace{4pt}}
\newcommand{\bx}[3]{{\sffamily\bfseries\fontsize{7.3}{9.2}\selectfont\color{#1!82!black}#2}\\[1.5pt]#3}
\newcommand{\cs}[2]{\expandafter\newcommand\csname CS@#1\endcsname{#2}}

% theorems
\newtheoremstyle{plainx}{7pt plus 2pt}{7pt plus 2pt}{\itshape}{}{\bfseries}{.}{0.5em}{}
\newtheoremstyle{uprightx}{7pt plus 2pt}{7pt plus 2pt}{}{}{\bfseries}{.}{0.5em}{}
\theoremstyle{plainx}
\newtheorem{theorem}{Theorem}[section]
\newtheorem{lemma}[theorem]{Lemma}
\newtheorem{proposition}[theorem]{Proposition}
\newtheorem{corollary}[theorem]{Corollary}
\theoremstyle{uprightx}
\newtheorem{definition}[theorem]{Definition}
\newtheorem{setting}[theorem]{Setting}          % hypotheses shared by several statements; each of them cites it by number
\newtheorem{remark}[theorem]{Remark}
\newtheorem{example}[theorem]{Example}
\renewcommand{\qedsymbol}{{\setlength{\fboxsep}{0pt}\setlength{\fboxrule}{0.5pt}\fbox{\rule{0pt}{0.46em}\hspace{0.46em}}}}
% what holds throughout the note: the ambient objects only (implicit and instance arguments), never a hypothesis
\newenvironment{conventions}{\par\addvspace{5pt}\noindent\textbf{Conventions.}\enspace\ignorespaces}{\par\addvspace{3pt}}
\makeatletter
\renewenvironment{proof}[1][Proof]{\par\pushQED{\qed}\normalfont\topsep0pt\partopsep0pt
  \trivlist\item[\hskip\labelsep\itshape #1.]\ignorespaces}{\popQED\endtrivlist\@endpefalse}
\makeatother

% where each declaration of the concordance is treated: \cs{key}{place}, the place naming a \label or a lettered item
% where each declaration of the concordance is treated: \cs{key}{place}, the place naming a \label or a lettered item
\cs{hasLogBoundvarpi}{Lemma~\ref{l:sizes}}
\cs{hasLogBoundrootDisc}{Lemma~\ref{l:sizes}}
\cs{piGBar}{Definition~\ref{d:chosen}}
\cs{IsChosen}{Definition~\ref{d:chosen}}
\cs{ChosenPlace}{Definition~\ref{d:chosen}}
\cs{chosen}{Definition~\ref{d:chosen}}
\cs{piGBarmulstar}{Lemma~\ref{l:chosen}}
\cs{piGnezero}{Lemma~\ref{l:chosen}}
\cs{piGBarnezero}{Lemma~\ref{l:chosen}}
\cs{existscofactor}{Lemma~\ref{l:chosen}}
\cs{isChoseniff}{Lemma~\ref{l:chosen}}
\cs{qeqg}{Lemma~\ref{l:chosen}}
\cs{finiteisChosen}{Lemma~\ref{l:chosen}}
\cs{existsresiduecharoforder}{Lemma~\ref{l:reschar}}
\cs{decla}{Definition~\ref{d:omega}}
\cs{orderOf}{Lemma~\ref{l:omega}}
\cs{omegaeqoneofnotisChosen}{Lemma~\ref{l:omega}}
\cs{omegazpowmulpoweqone}{Lemma~\ref{l:omega}}
\cs{existscharhasWeight}{Theorem B}
\cs{muFlat}{Definition~\ref{d:muflat}}
\cs{declb}{Definition~\ref{d:muflat}}
\cs{decld}{Definition~\ref{d:muflat}}
\cs{muFlathasWeight}{Lemma~\ref{l:muflat}}
\cs{chihasWeight}{Lemma~\ref{l:muflat}}
\cs{muhasWeight}{Theorem B}
\cs{Pair}{Definition~\ref{d:pair}}
\cs{logLevelTarget}{Definition~\ref{d:target}}
\cs{logLevelTargetmemIcc}{Theorem A}
\cs{N}{Definition~\ref{d:level}}
\cs{Nprime}{Theorem A}
\cs{NmodEqone}{Theorem A}
\cs{logNmemwindow}{Theorem A}
\cs{logNmemIcc}{Theorem A}
\cs{existspimulstar}{Theorem A}
\cs{declc}{Definition~\ref{d:pi}}
\cs{mulstar}{Theorem A}
\cs{TameChar}{Definition~\ref{d:tame}}
\cs{levelNorm}{Definition~\ref{d:tame}}
\cs{Twist}{Definition~\ref{d:tame}}
\cs{one}{Definition~\ref{d:tame}}
\cs{liftChar}{Definition~\ref{d:liftchar}}
\cs{liftTamePart}{Definition~\ref{d:liftchar}}
\cs{hasWeightliftChar}{Theorem B}
\cs{LiftIndex}{Definition~\ref{d:liftchar}}
\cs{BlockIndex}{Definition~\ref{d:block}}
\cs{firstBlock}{Definition~\ref{d:block}}
\cs{equivSum}{Definition~\ref{d:block}}
\cs{instFintypeYT}{Lemma~\ref{l:block}}
\cs{card}{Lemma~\ref{l:block}}
\cs{decle}{Definition~\ref{d:index}}
\cs{InitIndex}{Definition~\ref{d:index}}
\cs{X}{Definition~\ref{d:index}}
\cs{shape}{Definition~\ref{d:index}}
\cs{RowIndex}{Definition~\ref{d:index}}
\cs{Y}{Definition~\ref{d:index}}
\cs{cardrowIndex}{Lemma~\ref{l:rowindex}}
\cs{basisNumber}{Definition~\ref{d:rowpar}}
\cs{rowParameter}{Definition~\ref{d:rowpar}}
\cs{paramW}{Definition~\ref{d:rowpar}}
\cs{signsW}{Definition~\ref{d:rowpar}}
\cs{entrySign}{Definition~\ref{d:esign}}
\cs{iteoddeqsubemod}{Lemma~\ref{l:ite}}
\cs{G}{Definition~\ref{d:G}}
\cs{goodTupleFamilies}{Definition~\ref{d:G}}
\cs{e}{Definition~\ref{d:G}}
\cs{existsbranchingpairing}{Theorem C}
\cs{Node}{Definition~\ref{d:grid}}
\cs{NodeIndex}{Definition~\ref{d:grid}}
\cs{gridSpace}{Definition~\ref{d:grid}}
\cs{gridKernel}{Definition~\ref{d:gridforms}}
\cs{GridForm}{Definition~\ref{d:gridforms}}
\cs{gridLevelForms}{Definition~\ref{d:gridforms}}
\cs{gridChar}{Definition~\ref{d:gridchar}}
\cs{gridTamePart}{Definition~\ref{d:gridchar}}
\cs{triangleOrder}{Definition~\ref{d:gridlevels}}
\cs{nodeLevel}{Definition~\ref{d:gridlevels}}
\cs{levelIndex}{Definition~\ref{d:gridlevels}}
\cs{pointDescent}{Definition~\ref{d:moves}}
\cs{iotaPlus}{Definition~\ref{d:moves}}
\cs{descentMove}{Definition~\ref{d:moves}}
\cs{liftMove}{Definition~\ref{d:moves}}
\cs{isConjKernelPairgridKernel}{Lemma~\ref{l:gridkernel}}
\cs{gridCharhasWeight}{Lemma~\ref{l:gridcharweight}}
\cs{gridChardet}{Lemma~\ref{l:gridcharweight}}
\cs{isEdgeFamilygrid}{Lemma~\ref{l:edge}}
\cs{gridSpacediag}{Lemma~\ref{l:diag}}
\cs{gridSpacesubdiag}{Lemma~\ref{l:diag}}
\cs{levelIndexequntwisted}{Lemma~\ref{l:levelindex}}

% ------------------------------------------------------------------ macros of this note that are not notations of
% defined functions of the unit (those are made by the build, in data.tex, from the spec)
\newcommand{\OO}{\mathfrak{o}}
\newcommand{\CC}{\mathbb{C}}
\newcommand{\RR}{\mathbb{R}}
\newcommand{\ZZ}{\mathbb{Z}}
\newcommand{\QQ}{\mathbb{Q}}
\newcommand{\NN}{\mathbb{N}}
\newcommand{\FF}{\mathbb{F}}
\newcommand{\Ad}{\mathbb{A}}
\newcommand{\Af}{\mathbb{A}_f}
\newcommand{\Ohat}{\hat{\OO}}
\newcommand{\GL}{\mathrm{GL}}
\newcommand{\Un}{\mathrm{U}}
\newcommand{\cj}[1]{\overline{#1}}
\newcommand{\Nm}{\mathrm{N}}                       % the absolute norm of a finite place: \Nm w
\newcommand{\pt}{\ast}                             % the added point of an index set
% objects of other units, written with all their arguments (Chapter III: units 152; Chapter VII: units 140, 147)
\newcommand{\KP}[1]{K(#1)}
\newcommand{\cP}[1]{C(#1)}
\newcommand{\ED}[1]{E(#1)}
\newcommand{\None}[1]{N_1(#1)}
\newcommand{\fD}[1]{f(#1)}
\newcommand{\xD}[1]{x(#1)}
\newcommand{\xoD}[1]{x_1(#1)}
\newcommand{\bD}[1]{b(#1)}
\newcommand{\mD}[1]{m(#1)}
\newcommand{\rD}[1]{r(#1)}
\newcommand{\kD}[1]{k(#1)}
\newcommand{\tD}[1]{t(#1)}
\newcommand{\aD}[1]{a(#1)}
\newcommand{\gD}[1]{g(#1)}
\newcommand{\hD}[1]{h(#1)}
\newcommand{\hhD}[1]{h_2(#1)}
\newcommand{\nD}[2]{n_{#2}(#1)}
\newcommand{\wt}[3]{w_{#1,#2}(#3)}
\newcommand{\nw}[2]{w_{#1}(#2)}
\newcommand{\bone}[1]{b_1(#1)}
\newcommand{\btwo}[1]{b_2(#1)}
\newcommand{\thick}[1]{\operatorname{thick}(#1)}
\newcommand{\rdD}[1]{\operatorname{rd}_{\ED{#1}}}
\newcommand{\PlD}[1]{\operatorname{Pl}(#1)}
\newcommand{\Mom}[2]{M_\omega(#1,#2)}
\newcommand{\vpi}[2]{\varpi(#1,#2)}
\newcommand{\OmP}[2]{\Omega(#1,#2)}
\newcommand{\rtu}[1]{\varsigma_{#1}}               % the root of unity of order N_1(D) that generates E(D) over K
\newcommand{\ups}{\upsilon}                        % the splitting character on U(1)
\newcommand{\SpY}[2]{Y_{#2}(#1)}                   % the hermitian space Y_j(D) of unit 140
% the definitions of the unit whose census arguments are not Lean's (see Notes on the formal text)
\newcommand{\InitIdx}[2]{\underline{X}^{\circ}_{#2}(#1)}
\newcommand{\esign}[3]{\operatorname{sg}(#1,#2,#3)}
\newcommand{\nodeLev}[2]{\mathcal{K}_{D}(#1;#2)}
\newcommand{\oneT}{1_{D}}                          % the trivial tame character of D
% grids (unit 147)
\newcommand{\GrIdx}[3]{\underline{\mathcal{V}}^{#1}_{#2,#3}(D)}
\newcommand{\gsigns}[4]{\operatorname{sgn}^{#1,#2}_{#3,#4}(D)}
\newcommand{\gpar}[5]{\operatorname{gpar}^{#1,#2}_{#3,#4}(D,#5)}
\newcommand{\gwt}[3]{\operatorname{gw}_{D}(#1,#2,#3)}
\newcommand{\ghar}[9]{\operatorname{har}_{D}(#1,#2,#3,#4;#5;#6,#7;#8,#9)}
\newcommand{\gmod}[1]{\mathfrak{m}_{D}(#1)}
\newcommand{\tmod}[2]{\mathfrak{m}^{\triangle}_{D}(#1,#2)}
\newcommand{\lmod}[2]{\mathfrak{m}_{D}(#1;#2)}
% weights (background library)
\newcommand{\hc}[4]{\operatorname{hc}_{#1,#2}^{#3}(#4)}   % hc_{kappa,e}^{N}(tau,p): doubled Harish-Chandra parameter
\newcommand{\wrep}[2]{\rho_{#1,#2}}
\newcommand{\boxes}[2]{\operatorname{boxes}(#1,#2)}
\newcommand{\wn}[1]{\lVert #1\rVert}
\newcommand{\Fof}[3]{F^{#1,#2}_{#3}}
\newcommand{\clint}[2]{\operatorname{mean}_{#1}(#2)}
\newcommand{\Kmax}{\mathcal{K}_0}
\newcommand{\dg}[1]{\operatorname{diag}_{\sigma}(#1)}
\newcommand{\chow}[1]{\operatorname{cs}(#1)}                 % Euler's constant minus Re(L'/L)(1,chi_K)
\newcommand{\rtuz}{\varsigma_0}                     % a primitive root of unity of order N_0 in C
\newcommand{\brm}{\mathrm{br}}                      % the branching matrix
\newcommand{\kup}{\kappa^{\uparrow}}
\newcommand{\kdn}{\kappa^{\downarrow}}
\newcommand{\knat}{\kappa^{\natural}}
\newcommand{\enat}{e^{\natural}}
\newcommand{\imag}{\mathrm{i}}

\begin{document}
\color{ink}

{\raggedright
{\sffamily\bfseries\fontsize{7.2}{9}\selectfont\color{leancol}\addfontfeatures{LetterSpace=8}READBACK NOTE 24\par}
\vspace{5pt}
{\fontsize{21.5}{25.5}\selectfont\bfseries The characters, the level prime and the index sets of the data; the grid spaces\par}
\vspace{5pt}
{\fontsize{12.4}{16}\selectfont\itshape\color{inksoft} One unit of the Lean artifact \ArtifactName, read back into
mathematics\par}}
\vspace{8pt}
{\color{rulegrey}\hrule height 0.6pt}
\vspace{7pt}
\thispagestyle{empty}

\textbf{Summary.} Besides scales and ranks, the data~$D$ of the argument carry a list of chosen objects from which
the automorphic forms of Chapters V to VII are built. This note fixes them: the chosen places above the two primes
$g_0,g_1$ and the residue characters there; three characters of the norm-one group $\Un(1)$ with their weights and
tame parts; the level, a prime $\levelN{D}$ in a short window, with an integer of norm $\levelN{D}$; tame characters
and twists; the index sets of the hermitian spaces of the chain of theta lifts, the finite group $\Ggrp{D}$ and its
idempotent; and the spaces, characters and levels on the two grids of lifts of Chapter~VII. Three statements have
content: the level exists and $m\log b\le\log\levelN{D}\le30\,m\log b$ (Theorem~A); the characters exist, by Weil's
criterion for algebraic Hecke characters (Theorem~B); and for interlacing weights there is an integral pairing that
computes periods over a definite unitary group as finite means (Theorem~C).

\statement{leancol}{\fontsize{9.6}{13.2}\selectfont\textbf{What this note is.} The artifact \ArtifactName{} is a
body of Lean code whose final theorems state the Riemann hypothesis and its analogue for Dirichlet $L$-functions.
That is a claim: its authors report that Lean's kernel accepts the code; the record of checks of this write-up says
who has built what; nothing was compiled for this note, and no mathematician has read the proof through. This note is one of four on one unit, the folder of the seesaw identity (Chapter~V): it treats
\UnitInNote{} of the unit's \UnitMembers{} declarations, those of the file of index types and of the file of the
grid spaces, which are vocabulary for other units; the notes \emph{The seesaw identity for the chain of theta
lifts}, \emph{Two consecutive theta lifts} and \emph{Chains of definite theta lifts and the Gauss-sum operator}
treat the rest. Every statement below was compared, argument by argument, with the Lean text. Proofs are given in
full, except that results of other units and of the background library are quoted as \emph{Inputs}, each with its
owner. The Lean text read is \PinText; Appendix~A is the concordance of \ConcDecls{} declarations.}

\sect{The statements}{s:intro}
The notation is that of Section~\ref{s:setting}. In particular $D$ is admissible data for parameters $P=(N_0,d)$,
with scales $x,x_1,b$, exponent~$\mu$ and number of lifts~$m$; $\ED{D}$ is its CM field, with infinite
places~$\tau$; $\wt{D}{j}{\tau}$ are the weights of Chapter~III; and $\logTarget{D}$, $\levelN{D}$, $\piN{D}$,
$\omD{D}{w}$, $\muZero{D}{j}$, $\liftChar{D}{\vartheta}{j}$, $\liftTame{D}{\vartheta}{j}$ are defined in
Sections \ref{s:chosen} to~\ref{s:tame}.

\statement{arith}{\phantomsection\label{t:A}\textbf{Theorem A (the level).}\enspace In Setting~\ref{set:data}, for
parameters $P$ and admissible data $D$ for~$P$:
\begin{enumerate}[label=(\roman*)]
\item $m\log b\le\logTarget{D}\le29\,m\log b$;
\item $\levelN{D}$ is a prime number, $\levelN{D}\equiv1$ modulo $\None{D}\,d$, and
\[\logTarget{D}\le\log\levelN{D}\le\logTarget{D}+\sqrt{\log b},\qquad m\log b\le\log\levelN{D}\le30\,m\log b;\]
\item there is an integer $\pi$ of $\ED{D}$ with $\pi\,\cj{\pi}=\levelN{D}$, and
$\piN{D}\cdot\cj{\piN{D}}=\levelN{D}$.
\end{enumerate}}
\leanline{\lean{logLevelTargetmemIcc}\sep \lean{Nprime}\sep \lean{NmodEqone}\sep \lean{logNmemwindow}\sep
\lean{logNmemIcc}\sep \lean{existspimulstar}\sep \lean{mulstar}}

\statement{arith}{\phantomsection\label{t:B}\textbf{Theorem B (the characters).}\enspace In Setting~\ref{set:data},
for parameters $P$ and admissible data $D$ for~$P$:
\begin{enumerate}[label=(\roman*)]
\item for all natural numbers $c$ and~$n$ there is a character of $\Un(1)$ of $\ED{D}$ (a continuous character of
the norm-one adeles, trivial on the norm-one elements of $\ED{D}$) of weight \m{\tau\mapsto-c\cdot\wt{D}{0}{\tau}}
and tame part \m{(\omD{D}{w}^{\,n})_w};
\item for every natural number $j$, the character $\muZero{D}{j}$ has weight \m{\tau\mapsto-\wt{D}{j}{\tau}} and tame
part \m{(\omD{D}{w}^{\,j})_w};
\item for every tame character $\vartheta$ of~$D$ and every natural number~$j$, the character
$\liftChar{D}{\vartheta}{j}$ has weight \m{\tau\mapsto-\wt{D}{j}{\tau}} and tame part $\liftTame{D}{\vartheta}{j}$.
\end{enumerate}}
\leanline{\lean{existscharhasWeight}\sep \lean{muhasWeight}\sep \lean{hasWeightliftChar}}

Theorem~C does not involve the data. A \kw{weight} for a CM field $E$ is a pair $(\kappa,e)$: for every infinite
place $\tau$ of~$E$ a function \m{\kappa_\tau\colon\NN\to\NN} (the number $\kappa_\tau(h)$ of columns of height $h$
of a Young diagram) and an integer $e_\tau$. In rank~$N$ its \kw{doubled Harish-Chandra parameter} is
\[\hc{\kappa}{e}{N}{\tau,p}=2\Bigl(e_\tau+\sum_{p<h<N}\kappa_\tau(h)\Bigr)+N-1-2p\qquad(0\le p<N),\]
and a family $l=(l_\tau)$ of maps from an index set of $N$ elements to $\ZZ$ \kw{is the doubled Harish-Chandra
parameter of $(\kappa,e)$} if every $l_\tau$ is injective and its values, in decreasing order, are
$\hc{\kappa}{e}{N}{\tau,0},\dots,\hc{\kappa}{e}{N}{\tau,N-1}$. The representation $\wrep{\kappa}{e}$, the Weyl norm
$\wn{\cdot}$, the functions $\Fof{l}{\mathbf{s}}{\eta}$ and the mean $\clint{\mathcal{K}}{\cdot}$ are recalled in
Section~\ref{s:branching}.

\statement{arith}{\phantomsection\label{t:C}\textbf{Theorem C (branching pairing).}\enspace Let $E$ be a CM field
with a CM type, $n$ a non-empty finite index set and $n'$ a finite index set, $(\kappa,e)$ and $(\kappa',e')$
weights for~$E$, $l=(l_\tau)$ and $l'=(l'_\tau)$ families of maps $l_\tau\colon n\to\ZZ$, $l'_\tau\colon n'\to\ZZ$,
$\sigma\colon n\sqcup\{\pt\}\to n'$ a bijection, and $\mathbf{s}$, $\mathbf{r}$, $\mathbf{t}$ signs on $n$, on
$\{\pt\}$ and on~$n'$. Suppose that $(\mathbf{s}\sqcup\mathbf{r})\circ\sigma^{-1}=\mathbf{t}$ and $\mathbf{t}=1$;
that $l'$ is the doubled Harish-Chandra parameter of $(\kappa,e)$ and $-l$ is the
doubled Harish-Chandra parameter of $(\kappa',e')$; and that $l$ and $l'$ interlace: for every infinite
place~$\tau$ and every integer~$y$,
\[\#\{j\in n':y<l'_\tau(j)\}-\#\{i\in n:y<l_\tau(i)\}\in\{0,1\}.\]
Then there is a matrix $B$ with integer entries, from coordinates of weight $\kappa$ in rank $|n'|$ to coordinates of
weight $\kappa'$ in rank~$|n|$, with the following properties.
\begin{enumerate}[label=(\alph*)]
\item There are integers $N_\tau$, one for each infinite place~$\tau$, such that:
  \begin{enumerate}[label=(\arabic*)]
  \item for all commutative rings $S$ with an involution, all fields $R$ of characteristic~$0$, all families of ring
  homomorphisms $f_\tau\colon S\to R$, all $g\in\Un_n(S)$, $u\in\Un_1(S)$ and all vectors $v'$, $v$ over~$R$,
  \[\bigl\langle\wrep{\kappa'}{e'}\bigl(f(g)\bigr)v',\,B\,\wrep{\kappa}{e}\bigl(f(\dg{g,u})\bigr)v\bigr\rangle
  =\prod_\tau f_\tau(\det u)^{N_\tau}\cdot\langle v',Bv\rangle;\]
  \item $|N_\tau|\le\boxes{\kappa}{|n'|}+|e_\tau|$ for every~$\tau$;
  \item for every natural number $M\ge\boxes{\kappa}{|n'|}$ and all complex vectors $v'$, $v$:
  \m{|\langle v',Bv\rangle|\le(M^{M})^{3}\,\wn{v'}\,\wn{v}}.
  \end{enumerate}
\item For every point $g=(g_\tau)_\tau$ of the product, over the infinite places, of the real unitary groups of the
signs~$\mathbf{t}$ there is a complex number $C$ such that
for every standard hermitian space $V$ with index set $n$ and signs~$\mathbf{s}$, all algebraic forms $\eta$ of
weight $(\kappa',e')$ in rank~$|n|$ and $\phi$ of weight $(\kappa,e)$ in rank~$|n'|$, every $u\in\Un_1(\Af)$ and every
level $\mathcal{K}\subseteq\Un_n(\Af)$,
\[\begin{aligned}\int_{[\Un(V)]}\Fof{-l}{\mathbf{s}}{\eta}(y)\cdot{}&\Fof{l'}{\mathbf{t}}{\phi}\bigl(\dg{y,(1,u)}\cdot(g,1)\bigr)\,dy\\
&=\frac{C}{[\Kmax:\mathcal{K}]}\cdot\clint{\mathcal{K}}{h\mapsto\langle\eta(h),B\,\phi(\dg{h,u})\rangle}.\end{aligned}\]
\end{enumerate}}
\leanline{\lean{existsbranchingpairing}}

\sect{The setting}{s:setting}
\begin{conventions}\phantomsection\label{c:conventions}
Throughout, $P$ denotes parameters and $D$ admissible data for~$P$; Setting~\ref{set:data} recalls what these are,
and every statement that takes $D$ as an argument cites it. The letter $\tau$ denotes an infinite place of the
field $\ED{D}$ and $w$ a finite place of it; $i\in\{0,1\}$ numbers the two primes $g_0,g_1$ of~$D$. A stage is a
natural number~$j$; $\vartheta$ is a tame character of~$D$ (Definition~\ref{d:tame}), in Section~\ref{s:grid} an
$\mD{D}$-fold twist; $\psi$ is a family $(\psi_w)_w$ of characters of the residue fields of~$\ED{D}$. In
Definition~\ref{d:block} and Lemma~\ref{l:block}, $S$, $T$ and $\alpha$ are sets. In Theorem~C, $E$ is a CM field,
$n$ and $n'$ are finite index sets, $n$ non-empty, and $\kappa,\kappa'$, $e,e'$, $l,l'$,
$\mathbf{s},\mathbf{r},\mathbf{t}$, $\sigma$ are as stated there; in Lemma~\ref{l:gridcharweight}, $n$ and $n'$ are
finite index sets. In Section~\ref{s:grid}, $z$ and $c$ are elements of $\{0,1\}$, $I$ is a set, and
$i,i',j,j',a,l$ are natural numbers, and the offsets $u$, $v$ of a node on the grids (Definition~\ref{d:grid}) are natural numbers. Every other letter is bound by the statement that uses it.
\end{conventions}

\begin{setting}[parameters and data]\label{set:data}
\kw{Parameters} are a pair $P=(N_0,d)$: a positive multiple $N_0$ of $18480$ and a prime $d\equiv N_0/2-1$
modulo~$N_0$; put $\KP{P}=\QQ(\sqrt{-d})$. \kw{Admissible data} for $P$ are a record $D$ with these components: real
scales $\xD{D}$, $\xoD{D}$, $\bD{D}$; natural numbers $\mu$ (the exponent), $\mD{D}$ (the number of lifts) and $\rD{D}$ (the rank); a set
$\thick{D}$ of infinite places of the field $\ED{D}$ below; two primes $g_0\ne g_1$; and an integer $\pi_g$ of $\KP{P}$
with \m{\pi_g\cj{\pi_g}=g_0g_1}. They satisfy the conditions of Chapter~III. This note uses the following ones, in
which $\cP{P}\ge64$ is the constant of~$P$ of Chapter~III:
\begin{enumerate}[label=(D\arabic*)]
\item the scales satisfy $10^{12}\,\cP{P}\le\log\log \bD{D}$, $100\log\log \bD{D}\le\log \xD{D}$, $10^6\log \xD{D}\le\log \xoD{D}$ and
$10^6\log \xoD{D}\le\log \bD{D}$;
\item $\mD{D}$ is a power of $2$, $\mD{D}\ge2$, and $2^\mu=\mD{D}^{22}$;
$\sqrt{\xD{D}}\le \mD{D}\le2\sqrt{\xD{D}}$; $\mu\le32\log \xD{D}$; $\rD{D}$ is a power of $3$ with $\xD{D}^5\le \rD{D}\le3\xD{D}^5$;
\item for $i\in\{0,1\}$: \m{g_i\equiv1} modulo $\None{D}\,d$, the square of $\Mom{P}{\mu}$ divides $g_i-1$, and
\m{\bD{D}^{10}\le g_i\le \bD{D}^{11}};
\item for every real $U$ with \m{100\log \bD{D}\le U\le2^\mu} there is a prime $p\equiv1$ modulo $\Mom{P}{\mu}^2$ with
\m{U\le\log p\le U+\sqrt{\log \bD{D}}}.
\end{enumerate}
The following objects of Chapter~III are attached to $P$ and~$D$. The number
$\None{D}=N_0\cdot3^{\lfloor\log_3\mu\rfloor+5}\cdot2^{\mu}$ and the field $\ED{D}$, generated over $\KP{P}$ by a
root of unity $\rtu{D}$ of order $\None{D}$: a CM field, with complex conjugation $y\mapsto\cj{y}$ and ring of
integers $\OO_{\ED{D}}$. The set $\PlD{D}$ of its infinite places, of cardinality $\fD{D}$; each $\tau\in\PlD{D}$
comes with an embedding \m{\ED{D}\to\CC}, also written~$\tau$, and \m{|y|_\tau=|\tau(y)|}. The root
discriminant~$\rdD{D}$. The finite places of $\ED{D}$ are the non-zero prime ideals $w$ of $\OO_{\ED{D}}$; $\FF_w$ is
the residue field, $\Nm w$ its cardinality, $\cj{w}$ the conjugate place, and $w$ \kw{lies above} a prime $p$ if
$p\in w$. The numbers \m{\kD{D}=\rD{D}+1-\mD{D}}, \m{\tD{D}=5\rD{D}}, \m{\hhD{D}=\rD{D}^8}, \m{\hD{D}=2\,\hhD{D}} and
\m{\nD{D}{j}=\kD{D}+j\,(\tD{D}+1)}; the least multiple $\aD{D}$ of $\None{D}$ that is at least~$\xD{D}\,\bD{D}$, and the least
multiples $\bone{D}$, $\btwo{D}$ of $\None{D}$ that are at least $\aD{D}/\xD{D}$ and $\aD{D}/\xoD{D}$. The \kw{weights}
\[\wt{D}{j}{\tau}=\begin{cases}\bone{D}+j\,\hD{D}+\hhD{D}&\text{if }\tau\in\thick{D},\\
\btwo{D}+j\,\hD{D}+\hhD{D}&\text{otherwise,}\end{cases}\]
and the \kw{nominal weight} $\nw{D}{\tau}$, equal to $\aD{D}/\xD{D}$ if $\tau\in\thick{D}$ and to $\aD{D}/\xoD{D}$ otherwise.
The number \m{\gD{D}=g_0g_1}. Finally \m{\Mom{P}{\mu}=(N_0d)^{\mu+6}}; $\vpi{P}{\mu}$ is the purely imaginary generator
of the different of $\ED{D}$ of Chapter~III; and \m{\OmP{P}{\mu}=e^{\chow{d}}\cdot2\pi/\rdD{D}}, where $\chow{d}$ is
Euler's constant minus the real part of the logarithmic derivative at $1$ of the $L$-function of the quadratic
character of~$\KP{P}$.
\end{setting}

\statement{ana}{\textbf{Inputs from Chapter III} (unit~152, the folder of the data; the notes \emph{The parameters and
the CM field} and \emph{Admissible data}).\enspace In Setting~\ref{set:data}:

(I1) If $p$ is a prime with $p\equiv1$ modulo $\None{D}\,d$, every finite place of $\ED{D}$ above $p$ has absolute
norm~$p$.

(I2) Every root of unity of $\ED{D}$ is a power of $\rtu{D}$; \m{\prod_{\tau\in\PlD{D}}\tau(y)} is the norm of $y$
from $\ED{D}$ to~$\KP{P}$, for every $y\in\ED{D}$; the cyclotomic polynomial of order $\None{D}$ is irreducible
over~$\KP{P}$.

(I3) $\None{D}\,d$ divides $\Mom{P}{\mu}$.

(I4) For every $\tau\in\PlD{D}$: \m{|\vpi{P}{\mu}|_\tau=\sqrt d\cdot|\Psi'(\rtu{D})|_\tau}, where $\Psi$ is the
cyclotomic polynomial of order~$\None{D}$; and \m{\Psi(X)=\Psi_0(X^{\None{D}/N_0})}, where $\Psi_0$ is the cyclotomic
polynomial of order~$N_0$. Moreover \m{\prod_{\tau\in\PlD{D}}|\vpi{P}{\mu}|_\tau=\rdD{D}^{\,\fD{D}}}.

(I5) \m{\cP{P}\ge d}, \m{\cP{P}\ge N_0}, \m{\cP{P}\ge|\chow{d}|}, and
\m{\cP{P}\ge\bigl|\log|\Psi_0'(\rtuz)|\bigr|} for every primitive root of unity $\rtuz\in\CC$ of order~$N_0$.}

\subhead{Characters of $\Un(1)$.} These definitions belong to the background library. Let $E$ be a CM field with a
CM type, so that every infinite place $\tau$ has an embedding $\tau\colon E\to\CC$. The group $\Un(1)(\Ad_E)$
consists of the adeles $u$ of $E$ with $u\cj{u}=1$; it contains the group $\Un(1)(E)$ of the elements of $E$ of
norm one to the totally real subfield. A \kw{character of $\Un(1)$ of $E$} is a continuous homomorphism from
$\Un(1)(\Ad_E)$ to the unit circle that is trivial on $\Un(1)(E)$. Such a character $\nu$ has \kw{weight}
$\mathbf{e}=(\mathbf{e}(\tau))_\tau$, a family of integers, and \kw{tame part} $\psi=(\psi_w)_w$, a family of
multiplicative characters $\psi_w$ of the residue fields $\FF_w$, if $\psi_w$ is trivial for all but finitely many $w$
and
\[\nu(u)=\prod_\tau u_\tau^{\mathbf{e}(\tau)}\cdot\prod_w\psi_w(u_w\bmod w)\]
for every $u\in\Un(1)(\Ad_E)$ whose finite part is integral; here $u_\tau\in\CC$ is the component of $u$ at~$\tau$.

\statement{ana}{\textbf{Inputs from the background library.}\enspace Let $E$ be a CM field with a CM type.

(I6) \emph{Existence} (unit~104). Let $\mathbf{e}$ be a family of integers indexed by the infinite places and
$\psi$ a family of characters of the residue fields, trivial outside a finite set $\Sigma$ of finite places. If
\m{\prod_\tau\tau(u)^{\mathbf{e}(\tau)}\cdot\prod_{w\in\Sigma}\psi_w(u\bmod w)=1} for every integer $u$ of $E$ with
$u\cj{u}=1$, then there is a character of $\Un(1)$ of $E$ of weight $\mathbf{e}$ and tame part~$\psi$.
\emph{Status.} The statement is of the type of the existence of algebraic Hecke characters with prescribed
infinity type, for which the catalogue cites A.~Weil~[W] and A.~Raghuram~[R]. It records the verdict of its second
reading: the cited existence result, for characters of the full idele class group, has a different domain and a
different compatibility condition, and the existence theorem of~[R] prescribes the pure infinity type only. So
(I6) is not established as a special case of the cited results; the relation is unresolved, and (I6) is used here
as the library proves it.

(I7) \emph{Products} (unit~104). If $\nu$, $\nu'$ have weights $\mathbf{e}$, $\mathbf{e}'$ and tame parts $\psi$,
$\psi'$, then $\nu\nu'$ has weight $\mathbf{e}+\mathbf{e}'$ and tame part $\psi\psi'$, and $\nu^{y}$ has weight
$y\,\mathbf{e}$ and tame part $\psi^{y}$ for every integer~$y$.

(I8) \emph{The splitting character.} If $E$ carries Weil data, with splitting character $\ups$ restricted to
$\Un(1)(\Ad_E)$, then $\ups^{y}$ has the constant weight $y$ and trivial tame part, for every integer~$y$. The field
$\ED{D}$ of Setting~\ref{set:data} carries Weil data (Chapter~III).}

\sect{Chosen places and their residue characters}{s:chosen}
\begin{definition}[chosen places]\label{d:chosen}
In Setting~\ref{set:data}, let $P$ be parameters and $D$ admissible data for~$P$.
\begin{enumerate}[label=(\alph*)]
\item $\piGbar{D}$ is the complex conjugate of the image of $\pi_g$ in the ring of integers of $\ED{D}$.
\item A finite place $w$ of $\ED{D}$ is a \kw{chosen place of~$D$}, in symbols $\isChosen{D}{w}$, if
$\piGbar{D}\in w$; otherwise we write $\notChosen{D}{w}$.
\item For $i\in\{0,1\}$, $\ChPl{D}{i}$ is the set of the finite places $w$ of $\ED{D}$ with $\piGbar{D}\in w$ and
$g_i\in w$: the chosen places above~$g_i$.
\item For a pair $x=(i,w)$ with $i\in\{0,1\}$ and $w\in\ChPl{D}{i}$, the place $\chosenOf{D}{x}$ is~$w$.
\end{enumerate}
\end{definition}
\leanline{\lean{piGBar}\sep \lean{IsChosen}\sep \lean{ChosenPlace}\sep \lean{chosen}}

\begin{lemma}\label{l:chosen}
In Setting~\ref{set:data}, let $P$ be parameters, $D$ admissible data for~$P$, $i\in\{0,1\}$ and $w$ a finite place
of~$\ED{D}$.
\begin{enumerate}[label=(\roman*)]
\item \m{\piGbar{D}\cdot\cj{\piGbar{D}}=\gD{D}}; $\pi_g\ne0$ and $\piGbar{D}\ne0$.
\item There is a natural number $c$, prime to $g_i$, with \m{\piGbar{D}\cdot\cj{\piGbar{D}}=g_i\,c}.
\item $\isChosen{D}{w}$ holds if and only if, for some $i\in\{0,1\}$, both $\isChosen{D}{w}$ and $g_i\in w$ hold.
\item If $w$ lies above $g_i$, then $\Nm w=g_i$.
\item The set of the chosen places of $D$ is finite; hence so are $\ChPl{D}{0}$ and $\ChPl{D}{1}$.
\end{enumerate}
\end{lemma}
\begin{proof}
(i) The product is the image of \m{\cj{\pi_g}\,\pi_g=g_0g_1}, and $g_0g_1\ne0$. (ii) Take $c=g_{1-i}$; the primes
$g_0$ and $g_1$ are distinct. (iii) If $\piGbar{D}\in w$ then $g_0g_1\in w$ by~(i), and $w$ is a prime ideal.
(iv) This is (I1), by~(D3). (v) A non-zero integer of $\ED{D}$ lies in finitely many prime ideals.
\end{proof}
\leanline{\lean{piGBarmulstar}\sep \lean{piGnezero}\sep \lean{piGBarnezero}\sep \lean{existscofactor}\sep
\lean{isChoseniff}\sep \lean{qeqg}\sep \lean{finiteisChosen}}

\begin{lemma}\label{l:reschar}
In Setting~\ref{set:data}, let $P$ be parameters, $D$ admissible data for~$P$ and $w$ a finite place of $\ED{D}$
with $\isChosen{D}{w}$. Then there is a multiplicative character $T$ of the residue field $\FF_w$, with values
in~$\CC$, of order exactly~$\Mom{P}{\mu}^2$.
\end{lemma}
\begin{proof}
By Lemma~\ref{l:chosen}(iii) and~(iv), $\FF_w$ has $g_i$ elements for some~$i$, so $\FF_w^\times$ is cyclic of order
$g_i-1$, which $\Mom{P}{\mu}^2$ divides by~(D3). Map a generator to a primitive root of unity of order
$\Mom{P}{\mu}^2$ in~$\CC$.
\end{proof}
\leanline{\lean{existsresiduecharoforder}}

\begin{definition}[the residue characters]\label{d:omega}
In Setting~\ref{set:data}, let $D$ be admissible data and $w$ a finite place of~$\ED{D}$. If $w$ is a chosen place
of~$D$, then \m{\omD{D}{w}=T_w^{\,\Mom{P}{\mu}}}, where $T_w$ is a chosen character of $\FF_w$ of order
$\Mom{P}{\mu}^2$ (Lemma~\ref{l:reschar}); otherwise $\omD{D}{w}$ is the trivial character of~$\FF_w$. We write
$\omD{D}{\bullet}$ for the family $(\omD{D}{w})_w$.
\end{definition}
\leanline{\lean{decla}}

\begin{lemma}\label{l:omega}
In Setting~\ref{set:data}, let $P$ be parameters, $D$ admissible data for~$P$ and $w$ a finite place of~$\ED{D}$.
\begin{enumerate}[label=(\roman*)]
\item If $\isChosen{D}{w}$, then $\omD{D}{w}$ has order exactly~$\Mom{P}{\mu}$.
\item If $\notChosen{D}{w}$, then $\omD{D}{w}=1$.
\item Let $n$ be an integer and $\psi$ a family of characters of the residue fields with $\psi_w=1$. Then
\m{\bigl(\omD{D}{w}^{\,n}\,\psi_w\bigr)^{\Mom{P}{\mu}}=1}.
\end{enumerate}
\end{lemma}
\begin{proof}
(i) A power $T^{M}$ of an element $T$ of order $M^2$ has order~$M$. (ii) By definition. (iii) By (i) and (ii),
\m{\omD{D}{w}^{\,\Mom{P}{\mu}}=1} in both cases.
\end{proof}
\leanline{\lean{orderOf}\sep \lean{omegaeqoneofnotisChosen}\sep \lean{omegazpowmulpoweqone}}

\sect{The characters of the data}{s:chars}
\begin{proof}[Proof of Theorem~B(i)]
Put \m{\mathbf{e}(\tau)=-c\,\wt{D}{0}{\tau}} and \m{\psi_w=\omD{D}{w}^{\,n}}. By Lemma~\ref{l:omega}(ii) and
Lemma~\ref{l:chosen}(v), $\psi$ is trivial outside the finite set $\Sigma$ of the chosen places, so by~(I6) it
suffices to show
\begin{equation}\label{e:unit}
\prod_{\tau\in\PlD{D}}\tau(u)^{-c\,\wt{D}{0}{\tau}}\cdot\prod_{w\in\Sigma}\omD{D}{w}(u\bmod w)^{n}=1
\end{equation}
for every integer $u$ of $\ED{D}$ with $u\cj{u}=1$. Such a $u$ is an algebraic integer all of whose conjugates have
absolute value~$1$, since complex conjugation commutes with every embedding of a CM field; so $u$ is a root of
unity (Kronecker), and \m{u=\rtu{D}^{\,q}} for some natural number~$q$ by~(I2). By Setting~\ref{set:data},
\m{\wt{D}{0}{\tau}-\hhD{D}} is $\bone{D}$ or $\btwo{D}$, a multiple of~$\None{D}$; so
\m{\tau(u)^{\wt{D}{0}{\tau}}=\tau(u)^{\hhD{D}}}, and the first product in~\eqref{e:unit} is
\[\Bigl(\prod_{\tau\in\PlD{D}}\tau(\rtu{D})\Bigr)^{-c\,q\,\hhD{D}}.\]
By (I2) the product in brackets is the norm of $\rtu{D}$ to $\KP{P}$, and the minimal polynomial of $\rtu{D}$ over
$\KP{P}$ is the cyclotomic polynomial $\Psi$ of order~$\None{D}$. Its degree is even and $\Psi(0)=1$, because
$\None{D}>2$; so the norm is~$1$. For the second product, let $w\in\Sigma$. Then \m{\omD{D}{w}=T_w^{\,\Mom{P}{\mu}}}, and
\m{\rtu{D}^{\,\Mom{P}{\mu}}=1} because $\None{D}$ divides $\Mom{P}{\mu}$ by~(I3); hence
\m{\omD{D}{w}(\rtu{D}\bmod w)=T_w(\rtu{D}^{\,\Mom{P}{\mu}}\bmod w)=1}. This proves~\eqref{e:unit}.
\end{proof}

\begin{definition}[the three characters]\label{d:muflat}
In Setting~\ref{set:data}, let $D$ be admissible data.
\begin{enumerate}[label=(\alph*)]
\item $\muFlat{D}$ is a chosen character of $\Un(1)$ of $\ED{D}$ of weight \m{\tau\mapsto-\wt{D}{0}{\tau}} and
trivial tame part (Theorem~B(i) with $c=1$, $n=0$).
\item $\chiCirc{D}$ is a chosen character of $\Un(1)$ of $\ED{D}$ of weight $0$ and tame part $\omD{D}{\bullet}$
(Theorem~B(i) with $c=0$, $n=1$).
\item For a natural number~$j$,
\m{\muZero{D}{j}=\muFlat{D}\cdot\bigl(\ups^{-\hD{D}}\,\chiCirc{D}\bigr)^{j}}, where $\ups$ is the splitting character
of the Weil data of $\ED{D}$, restricted to~$\Un(1)$.
\end{enumerate}
\end{definition}
\leanline{\lean{muFlat}\sep \lean{declb}\sep \lean{decld}}

\begin{lemma}\label{l:muflat}
In Setting~\ref{set:data}, let $P$ be parameters and $D$ admissible data for~$P$. The character $\muFlat{D}$ has weight
\m{\tau\mapsto-\wt{D}{0}{\tau}} and trivial tame part, and $\chiCirc{D}$ has weight $0$ and tame part
\m{(\omD{D}{w})_w}.
\end{lemma}
\begin{proof}
By the choices of Definition~\ref{d:muflat}, since \m{\omD{D}{w}^{\,0}=1} and \m{\omD{D}{w}^{\,1}=\omD{D}{w}}.
\end{proof}
\leanline{\lean{muFlathasWeight}\sep \lean{chihasWeight}}

\begin{proof}[Proof of Theorem~B(ii) and (iii)]
(ii) By (I8) the character $\ups^{-\hD{D}}$ has weight $-\hD{D}$ and trivial tame part. By Lemma~\ref{l:muflat}
and~(I7), \m{\ups^{-\hD{D}}\chiCirc{D}} has weight $-\hD{D}$ and tame part $\omD{D}{\bullet}$, its $j$-th power has
weight $-j\,\hD{D}$ and tame part \m{(\omD{D}{w}^{\,j})_w}, and $\muZero{D}{j}$ has weight
\m{-\wt{D}{0}{\tau}-j\,\hD{D}=-\wt{D}{j}{\tau}} and the same tame part. (iii) By Definition~\ref{d:liftchar},
\m{\liftChar{D}{\vartheta}{j}=\muZero{D}{j}\cdot\vartheta} and $\vartheta$ has weight $0$ and tame
part~$\vartheta^{\mathrm{t}}$; apply (ii) and~(I7).
\end{proof}

\begin{definition}[pairs]\label{d:pair}
In Setting~\ref{set:data}, let $D$ be admissible data. The set $\PairSet{D}$ of the \kw{pairs} of~$D$ consists of the
quadruples $(\nu,\mathbf{c},\psi,z)$ of a character $\nu$ of $\Un(1)$ of~$\ED{D}$, a family
\m{\mathbf{c}=(\mathbf{c}(\tau))_{\tau\in\PlD{D}}} of natural numbers, a family $\psi$ of characters of the residue
fields and a natural number~$z$ such that: $\nu$ has weight \m{\tau\mapsto-\mathbf{c}(\tau)} and tame part~$\psi$;
$\psi_w=1$ for every finite place $w$ that is not a chosen place of~$D$; \m{1\le z\le\tD{D}}; and
\m{z\le\mathbf{c}(\tau)\le4\bD{D}} for every~$\tau$.
\end{definition}
\leanline{\lean{Pair}}

\sect{The level}{s:level}
\begin{lemma}[sizes]\label{l:sizes}
In Setting~\ref{set:data}, let $P$ be parameters and $D$ admissible data for~$P$. Then
\begin{enumerate}[label=(\roman*)]
\item \m{\xD{D}^{10}\le|\vpi{P}{\mu}|_\tau\le \xD{D}^{12}} for every infinite place $\tau$ of $\ED{D}$;
\item \m{\xD{D}^{10}\le\rdD{D}\le \xD{D}^{12}};
\item \m{\xD{D}^{-13}\le\OmP{P}{\mu}\le1}.
\end{enumerate}
\end{lemma}
\begin{proof}
(i) By (I4), \m{\Psi'(\rtu{D})=M_0\,\rtu{D}^{\,M_0-1}\,\Psi_0'(\rtu{D}^{\,M_0})} with
\m{M_0=\None{D}/N_0}, and $\rtu{D}^{\,M_0}$ is a primitive root of unity of order~$N_0$; so
\[\log|\vpi{P}{\mu}|_\tau=\tfrac12\log d+\log M_0+\log\bigl|\Psi_0'(\rtuz)\bigr|\]
for a primitive root of unity $\rtuz\in\CC$ of order~$N_0$. Here
\[\log M_0=(\lfloor\log_3\mu\rfloor+5)\log3+\mu\log2.\]
By (D2), \m{2^\mu=\mD{D}^{22}} lies between $\xD{D}^{11}$ and $2^{22}\xD{D}^{11}$, and
\[1\le3^{\lfloor\log_3\mu\rfloor+5}\le243\,\mu\le7776\log \xD{D}.\]
By (I5) the other two terms have absolute value at most $2\,\cP{P}$, and $10^{14}\,\cP{P}\le\log \xD{D}$ by~(D1).
Hence
\begin{multline*}11\log \xD{D}-2\,\cP{P}\le\log|\vpi{P}{\mu}|_\tau\\\le11\log \xD{D}+22\log2+\log(7776\log \xD{D})+2\,\cP{P},\end{multline*}
which lies between $10\log \xD{D}$ and $12\log \xD{D}$. (ii) By (I4), $\rdD{D}$ is the geometric mean of the numbers
$|\vpi{P}{\mu}|_\tau$; apply~(i).
(iii) By (I5), \m{e^{-\cP{P}}\le e^{\chow{d}}\le e^{\cP{P}}}; apply (ii) and~(D1).
\end{proof}
\leanline{\lean{hasLogBoundvarpi}\sep \lean{hasLogBoundrootDisc}}
In the formal text these bounds are \emph{instances}: facts that Lean's search for comparisons of sizes uses without
being named. (iii) is such an instance without a name.

\begin{definition}[the target]\label{d:target}
In Setting~\ref{set:data}, let $D$ be admissible data. Then
\[\logTarget{D}=\mD{D}\cdot\frac{1}{\fD{D}}\sum_{\tau\in\PlD{D}}\log\nw{D}{\tau}
+(\mD{D}-1)\cdot\log\frac{\gD{D}\cdot\rdD{D}}{2\pi}-\frac{15}{2}\log \bD{D} .\]
\end{definition}
\leanline{\lean{logLevelTarget}}

\begin{proof}[Proof of Theorem~A(i)]
By Setting~\ref{set:data}, $\xD{D}\,\bD{D}\le\aD{D}<\xD{D}\,\bD{D}+\None{D}$, and
$\None{D}\le\cP{P}\cdot243\,\mu\cdot2^{22}\xD{D}^{11}\le \bD{D}$ by (I5), (D2) and~(D1). So $\bD{D}\le\aD{D}/\xD{D}\le2\bD{D}$ and
$1\le \xD{D}\,\bD{D}/\xoD{D}\le\aD{D}/\xoD{D}\le2\bD{D}$, whence $0\le\log\nw{D}{\tau}\le2\log \bD{D}$ for every~$\tau$. By (D3),
$\bD{D}^{20}\le\gD{D}\le \bD{D}^{22}$; with Lemma~\ref{l:sizes}(ii) and $2\pi\le \xD{D}^{10}$, $\xD{D}^{12}\le \bD{D}^{5}$, this gives
$\bD{D}^{2}\le\gD{D}\,\rdD{D}/(2\pi)\le \bD{D}^{27}$. Therefore
\[(2\mD{D}-\tfrac{19}{2})\log \bD{D}\ \le\ \logTarget{D}\ \le\ \bigl(2\mD{D}+27(\mD{D}-1)\bigr)\log \bD{D} ,\]
and \m{\mD{D}\ge10} by (D2) and~(D1).
\end{proof}

\begin{definition}[the level]\label{d:level}
In Setting~\ref{set:data}, let $D$ be admissible data. By Theorem~A(i), (D1) and~(D2),
$100\log \bD{D}\le\logTarget{D}\le29\,\mD{D}\log \bD{D}\le \mD{D}^{22}=2^\mu$. The \kw{level} $\levelN{D}$ is a chosen prime number
with \m{\levelN{D}\equiv1} modulo $\Mom{P}{\mu}^2$ and
$\logTarget{D}\le\log\levelN{D}\le\logTarget{D}+\sqrt{\log \bD{D}}$; one exists by~(D4), applied to $U=\logTarget{D}$.
\end{definition}
\leanline{\lean{N}}

\begin{proof}[Proof of Theorem~A(ii) and (iii)]
(ii) $\levelN{D}$ is prime and lies in the window by its choice. By (I3), $\None{D}\,d$ divides $\Mom{P}{\mu}$, hence
$\Mom{P}{\mu}^2$, so \m{\levelN{D}\equiv1} modulo $\None{D}\,d$. The last two inequalities follow from (i) and
$\sqrt{\log \bD{D}}\le\log \bD{D}\le \mD{D}\log \bD{D}$.

(iii) Since $16$ divides $N_0$ and $N_0$ divides $\None{D}$, (ii) gives \m{\levelN{D}\equiv1} modulo~$4$. By Fermat's
theorem on sums of two squares there are integers $y$, $y'$ with \m{y^2+y'^2=\levelN{D}}. The element
\m{\imag=\rtu{D}^{\,\None{D}/4}} of $\OO_{\ED{D}}$ has $\imag^2=-1$, and \m{\cj{\imag}=\imag^{-1}=-\imag}
because complex conjugation inverts roots of unity. So \m{\pi=y+y'\,\imag} has
\m{\pi\,\cj{\pi}=y^2+y'^2=\levelN{D}}. The second statement is the choice of Definition~\ref{d:pi}.
\end{proof}

\begin{definition}\label{d:pi}
In Setting~\ref{set:data}, let $D$ be admissible data. Then $\piN{D}$ is a chosen integer $\pi$ of $\ED{D}$ with
\m{\pi\,\cj{\pi}=\levelN{D}} (Theorem~A(iii)).
\end{definition}
\leanline{\lean{declc}}

\begin{remark}\label{r:levelsizes}
Three more facts are recorded in the formal text as instances without names: $\levelN{D}$ is prime;
\m{\sqrt{\xD{D}}\,\log \bD{D}\le\log\levelN{D}\le e^{32}\sqrt{\xD{D}}\,\log \bD{D}}; and \m{\levelN{D}\ge \bD{D}^{101}}. They follow from
Theorem~A(ii) and \m{\sqrt{\xD{D}}\le \mD{D}\le2\sqrt{\xD{D}}}, \m{\mD{D}\ge101}. By (I1) and Theorem~A(ii), every finite place of
$\ED{D}$ above $\levelN{D}$ has absolute norm~$\levelN{D}$.
\end{remark}

\sect{Tame characters and twists}{s:tame}
\begin{definition}[tame characters, twists]\label{d:tame}
In Setting~\ref{set:data}, let $P$ be parameters and $D$ admissible data for~$P$.
\begin{enumerate}[label=(\alph*)]
\item The set $\TameSet{D}$ of the \kw{tame characters of~$D$} consists of the triples of a character of $\Un(1)$ of
$\ED{D}$, a finite set $Q$ of finite places of $\ED{D}$ and a family $\psi$ of characters of the residue fields, such
that: the character has weight $0$ and tame part~$\psi$; $\psi_w=1$ for every $w\notin Q$; $\cj{w}\notin Q$ for
every $w\in Q$; and \m{\gD{D}\cdot\levelN{D}} is prime to \m{\prod_{w\in Q}\Nm w}. We write $\vartheta$ for a tame
character and for its character of $\Un(1)$, $Q(\vartheta)$ for its set and
\m{\vartheta^{\mathrm{t}}=(\vartheta^{\mathrm{t}}_w)_w} for its family.
\item For a tame character $\vartheta$ of $D$, \m{\levelNorm{\vartheta}=\prod_{w\in Q(\vartheta)}\Nm w}.
\item The set $\TwistSet{D}$ of the \kw{twists of~$D$} consists of the tame characters $\vartheta\in\TameSet{D}$ with
\m{\vartheta^{\mathrm{t}}_w\ne1} for every $w\in Q(\vartheta)$.
\item The \kw{trivial tame character} $\oneT$ of $D$ is the trivial character with $Q(\oneT)=\emptyset$ and trivial
family; it is a twist of~$D$.
\end{enumerate}
\end{definition}
\leanline{\lean{TameChar}\sep \lean{levelNorm}\sep \lean{Twist}\sep \lean{one}}
A tame character may mark a place: its set $Q(\vartheta)$ may contain places $w$ with
\m{\vartheta^{\mathrm{t}}_w=1}. Non-triviality at every marked place is what a twist adds; so for a twist, and only then,
$Q(\vartheta)$ is exactly the set of the places where the tame part is not trivial. The formal text
makes the trivial twist the unit element $1$ of the type of twists and lets a twist be used as a tame character
(two instances without names). The $\mD{D}$-fold twists of the overview are the twists $\vartheta$ of $D$ such that
\m{(\vartheta^{\mathrm{t}}_w)^{\,n}\ne1} for all $w\in Q(\vartheta)$ and \m{1\le n\le \mD{D}}; they, and the power
$\vartheta^{n}$ of a tame character (the same set $Q(\vartheta)$, the family $(\vartheta^{\mathrm{t}})^{\,n}$), are defined in
Chapter~VII (unit~147).

\begin{definition}[characters of the lifts]\label{d:liftchar}
In Setting~\ref{set:data}, let $D$ be admissible data, $\vartheta$ a tame character of~$D$ and $j$ a natural number.
\begin{enumerate}[label=(\alph*)]
\item \m{\liftChar{D}{\vartheta}{j}=\muZero{D}{j}\cdot\vartheta}, a character of $\Un(1)$ of~$\ED{D}$.
\item $\liftTame{D}{\vartheta}{j}$ is the family \m{\bigl(\omD{D}{w}^{\,j}\cdot\vartheta^{\mathrm{t}}_w\bigr)_w} of characters
of the residue fields.
\item $\LiftIdx{D}$ is the set \m{\{1,\dots,\mD{D}\}}, the indices of the lifts.
\end{enumerate}
\end{definition}
\leanline{\lean{liftChar}\sep \lean{liftTamePart}\sep \lean{LiftIndex}}

\sect{Index sets}{s:index}
For a natural number $q$ we write $[q]=\{0,\dots,q-1\}$, and $\{\pt\}$ for a set with one element.

\begin{definition}[blocks]\label{d:block}
Let $S$, $T$ be sets and $j$ a natural number.
\begin{enumerate}[label=(\alph*)]
\item \m{\BlockIdx{j}{S}{T}} is defined by recursion: \m{\BlockIdx{0}{S}{T}=S} and
\m{\BlockIdx{j+1}{S}{T}=\BlockIdx{j}{S}{T}\sqcup T}. It is the disjoint union of $S$ and $j$ copies of~$T$.
\item \m{\firstBlock{j}\colon S\to\BlockIdx{j}{S}{T}} is the inclusion of the first block.
\item For a set $\alpha$ and a bijection $\nu\colon S\to\alpha$,
\m{\equivSum{\nu}{j}\colon\BlockIdx{j}{S}{T}\to\alpha\sqcup([j]\times T)} is a bijection: $\nu$ for $j=0$, and for
$j+1$ a chosen bijection, which exists because both sides have the cardinality of \m{\alpha\sqcup([j]\times T)\sqcup T}.
\end{enumerate}
\end{definition}
\leanline{\lean{BlockIndex}\sep \lean{firstBlock}\sep \lean{equivSum}}

\begin{lemma}\label{l:block}
Let $S$, $T$ be finite sets and $j$ a natural number. Then $\BlockIdx{j}{S}{T}$ is finite, with
\m{|S|+j\,|T|} elements.
\end{lemma}
\begin{proof}
By induction on~$j$; or by the bijection $\equivSum{\nu}{j}$ for the identity $\nu$ of~$S$.
\end{proof}
\leanline{\lean{instFintypeYT}\sep \lean{card}}

\begin{definition}[the index sets of the data]\label{d:index}
In Setting~\ref{set:data}, let $D$ be admissible data and $j$ a natural number.
\begin{enumerate}[label=(\alph*)]
\item \m{\rhoD{D}=\lfloor \mD{D}/2\rfloor-1}, or $0$ if this is negative.
\item \m{\InitIdx{D}{j}=[\nD{D}{j}-1]}.
\item \m{\Xidx{D}{j}=\InitIdx{D}{j}\sqcup\{\pt\}}: the basis index set of the standard definite hermitian space
$X_j(D)$ of rank~$\nD{D}{j}$, the point $\pt$ indexing the last basis vector.
\item $\shapeD{D}$ is the \kw{shape} consisting of the number $\rhoD{D}$, the index set $\Xidx{D}{0}$ and its
element~$\pt$.
\item \m{\Ridx{D}=[\rhoD{D}]\sqcup\bigl(\Xidx{D}{0}\sqcup\{\diamond\}\bigr)\sqcup[\rhoD{D}]}, the index set of the
shape $\shapeD{D}$: indices of $\rhoD{D}$ isotropic vectors, of $\Xidx{D}{0}$, of one further vector, and of
$\rhoD{D}$ dual isotropic vectors. The element $\diamond$ indexes the further vector, the middle vector of the
shape; it is not the point $\pt$, which is an element of~$\Xidx{D}{0}$.
\item \m{\Yidx{D}{j}=\BlockIdx{j}{\Ridx{D}}{[\tD{D}]}}: the basis index set of the hermitian space $\SpY{D}{j}$.
\end{enumerate}
\end{definition}
\leanline{\lean{decle}\sep \lean{InitIndex}\sep \lean{X}\sep \lean{shape}\sep \lean{RowIndex}\sep \lean{Y}}

\begin{lemma}\label{l:rowindex}
In Setting~\ref{set:data}, let $P$ be parameters and $D$ admissible data for~$P$. Then \m{\mD{D}=2\,\rhoD{D}+2}, the set
$\Ridx{D}$ has exactly $\rD{D}$ elements, and $\Yidx{D}{j}$ has \m{\rD{D}+j\,\tD{D}} elements for every natural number~$j$.
\end{lemma}
\begin{proof}
By (D2), $\mD{D}$ is even and at least~$2$, so $\rhoD{D}=\mD{D}/2-1$; and $\mD{D}\le2\sqrt{\xD{D}}\le \xD{D}^5\le \rD{D}$, so
\m{\kD{D}=\rD{D}+1-\mD{D}\ge1} and $\Xidx{D}{0}$ has \m{\nD{D}{0}=\kD{D}} elements. Hence
\[|\Ridx{D}|=2\,\rhoD{D}+\kD{D}+1=(\mD{D}-2)+(\rD{D}+1-\mD{D})+1=\rD{D}.\]
The last statement is Lemma~\ref{l:block}.
\end{proof}
\leanline{\lean{cardrowIndex}}
The formal text adds an instance without a name: a bound for the size of $\rD{D}$ is also one for~$|\Ridx{D}|$. The
identity \m{\mD{D}=2\,\rhoD{D}+2} is stated there in Chapter~VII (unit~140).

\begin{definition}[numbering, row parameter, signs]\label{d:rowpar}
In Setting~\ref{set:data}, let $D$ be admissible data.
\begin{enumerate}[label=(\alph*)]
\item \m{\basisNum{D}\colon\Ridx{D}\to[\rD{D}]} is a chosen bijection (Lemma~\ref{l:rowindex}).
\item For an infinite place $\tau$ of $\ED{D}$ and a natural number $i$, the entry $i$ of the \kw{row parameter} is
\[\rowPar{D}{\tau}{i}=\begin{cases}\kD{D}/2-1-i&\text{if }i+1<\kD{D},\\
\hhD{D}-\wt{D}{i+2-\kD{D}}{\tau}&\text{if }i+1\ge\kD{D}.\end{cases}\]
\item For an infinite place $\tau$ and $i\in\Ridx{D}$: \m{\parW{D}{\tau}{i}=2\,\rowPar{D}{\tau}{\basisNum{D}(i)}}.
\item For a natural number $j$ and $i\in\Ridx{D}$: \m{\signsW{D}{j}{i}=-1} if \m{\basisNum{D}(i)\ge\kD{D}+j-1}, and
\m{\signsW{D}{j}{i}=+1} otherwise.
\end{enumerate}
\end{definition}
\leanline{\lean{basisNumber}\sep \lean{rowParameter}\sep \lean{paramW}\sep \lean{signsW}}
The number $\kD{D}$ is even, since $\rD{D}$ is odd and $\mD{D}$ even. For $i=\kD{D}-1,\dots,\rD{D}-1$ the second case gives
\m{-(\wt{D}{q}{\tau}-\hhD{D})} with $q=i+2-\kD{D}\in\{1,\dots,\mD{D}\}$. So the row parameter consists of the
$\kD{D}-1$ numbers \m{\kD{D}/2-1,\dots,1-\kD{D}/2}, followed by $\mD{D}$ entries given by the weights of the $\mD{D}$ lifts; $\parW{D}{\tau}{\cdot}$ is the doubled parameter of the forms on $\SpY{D}{0}$ and on the
spaces $W_j$ of signs $\signsW{D}{j}{\cdot}$ of Chapter~VII.

\begin{definition}[entry signs]\label{d:esign}
For natural numbers $o$, $j$, $a$ the sign \m{\esign{o}{j}{a}\in\{\pm1\}} is defined by recursion on~$o$:
\[\begin{gathered}\esign{o+1}{0}{a}=-\esign{o}{0}{a},\\
\esign{o}{j}{a}=(-1)^{\,o+a+\min(\max(a-o,0),\;j)}\ \text{ in all other cases.}\end{gathered}\]
The second formula holds in the first case as well, since \m{\min(\max(a-o,0),0)=0}.
\end{definition}
\leanline{\lean{entrySign}}

\begin{lemma}\label{l:ite}
Let $j$ be a natural number and $y$ an integer. Then the number that is $y-1$ if $j$ is odd and $y$ if $j$ is even
equals \m{y-(j\bmod2)}.
\end{lemma}
\begin{proof}
$j\bmod2$ is $1$ or $0$ accordingly.
\end{proof}
\leanline{\lean{iteoddeqsubemod}}

\begin{definition}[the finite group and its idempotent]\label{d:G}
In Setting~\ref{set:data}, let $D$ be admissible data.
\begin{enumerate}[label=(\alph*)]
\item \m{\Ggrp{D}=\prod_{x}\GL_{\Ridx{D}}\bigl(\FF_{\chosenOf{D}{x}}\bigr)}, the product over the pairs $x=(i,w)$
with $i\in\{0,1\}$ and $w\in\ChPl{D}{i}$: a finite group, the product of the groups $\GL_r$ of the residue fields
of the chosen places.
\item $\goodFam{D}$ is the finite set of the families \m{(c_x)_x} of monic polynomials \m{c_x} of degree $r$ over
$\FF_{\chosenOf{D}{x}}$ such that, for each $i\in\{0,1\}$, the tuple \m{(c_{(i,w)})_{w\in\ChPl{D}{i}}} is \kw{good}
for $\Mom{P}{\mu}$ and~$\rD{D}$: every $c_{(i,w)}$ is squarefree with non-zero constant term and has no root $y$ in
$\FF_w$ with \m{y^{\Mom{P}{\mu}}=1}, and at least one $c_{(i,w)}$ is irreducible with a root of multiplicative order
\m{(\Nm w)^r-1} in an algebraic closure of~$\FF_w$.
\item \m{\ezero{D}=\sum_{c\in\goodFam{D}}\bigotimes_x e_{c_x}} in the group algebra $\CC[\Ggrp{D}]$, where $e_{c_x}$ is
the idempotent of Green's character with the label $c_x$ in the group algebra of
$\GL_{\Ridx{D}}(\FF_{\chosenOf{D}{x}})$.
\end{enumerate}
\end{definition}
\leanline{\lean{G}\sep \lean{goodTupleFamilies}\sep \lean{e}}
Good tuples and the idempotents $e_c$ belong to Chapter~VII (unit~140, the note \emph{Primitive labels of general
linear groups over finite fields}) and to the background library; there $\ezero{D}$ is shown to be a central
idempotent, and the rows of the overview are its image on the forms of principal level on~$\SpY{D}{0}$.

\sect{The branching pairing}{s:branching}
In this section $E$ is a CM field with a CM type, and weights and their Harish-Chandra parameters are as before
Theorem~C. We recall the objects of the background library that the theorem names.

\subhead{Representations.} For a weight $(\kappa,e)$ and a finite index set $n$, the vectors \kw{of weight $\kappa$
in rank~$|n|$} over a commutative ring $R$ are the coordinate vectors, in the basis of semistandard fillings, of the
product over the infinite places $\tau$ of the Weyl modules of $\GL_n$ whose Young diagram has $\kappa_\tau(h)$ columns
of height~$h$, for $0<h<|n|$. On them a family \m{g=(g_\tau)_\tau} of matrices in $\GL_n(R)$ acts by
\m{\wrep{\kappa}{e}(g)=\prod_\tau\det(g_\tau)^{e_\tau}} times the Schur matrix of~$g$. For vectors $v'$, $v$ with the
same index set, \m{\langle v',v\rangle=\sum_a v'_a v_a} is the bilinear dot product, and $\wn{v}^2$ is the value at
$(v,v)$ of the positive definite hermitian form on complex vectors of weight $\kappa$ given by the integral Gram
matrix of the Weyl lattice. \m{\boxes{\kappa}{N}=\sum_\tau\sum_{h<N}h\,\kappa_\tau(h)} is the number of boxes. For a
ring $S$ with an involution $y\mapsto\cj{y}$, $\Un_n(S)$ is the group of the $g\in\GL_n(S)$ with
\m{\cj{g}^{\mathrm{t}}g=1}; for $g\in\Un_n(S)$ and $u\in\Un_1(S)$, \m{\dg{g,u}\in\Un_{n'}(S)} is the block diagonal
matrix with blocks $g$ and~$u$, relabelled by~$\sigma$; and for ring homomorphisms $f_\tau\colon S\to R$, $f(g)$ is the
family \m{(f_\tau(g))_\tau}, entry by entry.

\subhead{Forms.} A \kw{standard} hermitian space $V$ with index set $n$ has Gram matrix $1$ and the standard
lattice; its signs are all~$+1$. Its adelic unitary group is the product of a compact real group and $\Un_n(\Af)$;
$[\Un(V)]$ is the quotient by the rational points, with the invariant measure that the background library fixes. An \kw{algebraic form} of
weight $(\kappa,e)$ in rank $|n|$ is a function $\phi$ on $\Un_n(\Af)$, invariant on the right under an open
subgroup, with values in the complex vectors of weight $\kappa$ in rank~$|n|$, such that
\m{\phi(y_0h)=\wrep{\kappa}{e}(y_0)\phi(h)} for rational points~$y_0$. For a family $l$ that is the doubled
Harish-Chandra parameter of $(\kappa,e)$ and signs~$\mathbf{s}$, the function $\Fof{l}{\mathbf{s}}{\phi}$ on the adelic
group is \m{(y_\infty,h)\mapsto\bigl(v^{+}_{l},\cj{\wrep{\kappa}{e}}(y_\infty)^{-1}\phi(h)\bigr)}, the coefficient of
the highest weight vector $v^{+}_{l}$ for the Gram form, where $\cj{\wrep{\kappa}{e}}$ is $\wrep{\kappa}{e}$ at the
complex conjugates of the matrices. A \kw{level} is an open subgroup $\mathcal{K}$ of the stabiliser $\Kmax$ of the
standard lattice in $\Un_n(\Af)$, and for a function $F$ on $\Un_n(\Af)$, left invariant under the rational points
and right invariant under an open subgroup, $\clint{\mathcal{K}}{F}$ is its integral over the quotient of
$\Un_n(\Af)$ by the rational points for the measure that gives $\mathcal{K}$ the volume~$1$: a finite sum over
classes, weighted by inverse orders of stabilisers.

\statement{ana}{\textbf{Inputs for Theorem C.}\enspace Let $E$, $n$, $n'$, $\sigma$ be as in Theorem~C, with
$|n'|=|n|+1$.

(I9) \emph{Parameters with a gap} (unit~143). If a family $l=(l_\tau)$ of maps $n\to\ZZ$ has, for every~$\tau$,
distinct values that are all congruent to $|n|-1$ modulo~$2$, then $l$ is the doubled Harish-Chandra parameter of
some weight $(\knat,\enat)$.

(I10) \emph{Dual pairing} (unit~101). If $l$ is the doubled Harish-Chandra parameter of
$(\knat,\enat)$ and $-l$ that of $(\kappa',e')$, both in rank~$|n|$, there is an integer matrix $A$
with: \m{\langle\wrep{\kappa'}{e'}(g)v',A\,\wrep{\knat}{\enat}(g)v\rangle=\langle v',Av\rangle} for
all fields $R$ of characteristic~$0$, all families $g$ in $\GL_n(R)$ and all $v'$, $v$; for all signs $\mathbf{s}$
and algebraic forms $\eta$ of weight $(\kappa',e')$,
\m{\Fof{-l}{\mathbf{s}}{\eta}(y_\infty,h)=\langle\eta(h),A\,\cj{\wrep{\knat}{\enat}}(y_\infty)v^{+}_{l}\rangle};
and for every complex vector $v$ of weight $\knat$ there is a complex vector $v^{\circ}$ of weight
$\kappa'$ with \m{\wn{v^{\circ}}=\wn{v}} and \m{\langle v',Av\rangle=(v^{\circ},v')} for all~$v'$, where
$(\cdot,\cdot)$ is the Gram form, conjugate linear in the first variable.

(I11) \emph{Branching} (unit~19; the branching law from $\GL_{|n|+1}$ to $\GL_{|n|}$ with multiplicity one, in an
integral form, here in the case $q=0$ of the formal statement; reference to be confirmed). Let $(\kappa,e)$ be a weight, and let $\kup$, $\kdn$ be families of
functions $\NN\to\NN$ with \m{\kup+\kdn=\kappa}. Let \m{(\knat,\enat)} be a weight with
\m{\knat_\tau(h)=\kdn_\tau(h)+\kup_\tau(h+1)} for $0<h<|n|$ and
\m{\enat_\tau=e_\tau+\kdn_\tau(|n|)}. Then there is an integer matrix $\brm$, from weight $\kappa$ in rank
$|n'|$ to weight $\knat$ in rank~$|n|$, such that: over every field $R$ of characteristic~$0$,
\[\brm\cdot\wrep{\kappa}{e}(\dg{g,u})=\prod_\tau u_\tau^{\,e_\tau+\sum_{0<h<|n'|}\kup_\tau(h)}\cdot\wrep{\knat}{\enat}(g)\cdot\brm\]
for all families $g$ in $\GL_n(R)$ and $u$ in $R^\times$; \m{\wn{\brm\,v}\le(M^{M})^{3}\wn{v}} for all complex $v$
and every natural number \m{M\ge\boxes{\knat}{|n|}}; and every complex matrix that intertwines
\m{\wrep{\kappa}{e}(\dg{g,1})} with \m{\wrep{\knat}{\enat}(g)} for all families $g$ of unitary matrices is a multiple
of~$\brm$.

(I12) \emph{Periods as means} (unit~27; reference to be confirmed). Let $V$ be a standard space with index set~$n$,
let $\brm$ have the last property of~(I11), let $L'$ be a linear functional on the complex vectors of weight
$\kappa$ in rank~$|n'|$ and $v^{\natural}$ a complex vector of weight $\knat$ in rank~$|n|$. Then there is
a complex number $C$ such that for every algebraic form $\phi$ of weight $(\kappa,e)$ in rank~$|n'|$, every automorphic
function $\Phi$ on the adelic unitary group of~$V$, every $u\in\Un_1(\Af)$, every family $(L_h)_h$ of linear
functionals on the complex vectors of weight $\knat$, indexed by $h\in\Un_n(\Af)$, and every level~$\mathcal{K}$: if
\m{\Phi(y_\infty,h)=L_h\bigl(\cj{\wrep{\knat}{\enat}}(y_\infty)v^{\natural}\bigr)} for all
$(y_\infty,h)$, then
\[\begin{aligned}\int_{[\Un(V)]}\Phi(y)\cdot{}&L'\bigl(\cj{\wrep{\kappa}{e}}(\dg{y_\infty,1})^{-1}\phi(\dg{h,u})\bigr)\,dy\\
&=\frac{C}{[\Kmax:\mathcal{K}]}\cdot\clint{\mathcal{K}}{h\mapsto L_h\bigl(\brm\,\phi(\dg{h,u})\bigr)},\end{aligned}\]
where $y=(y_\infty,h)$. The functions $\Fof{-l}{\mathbf{s}}{\eta}$ are automorphic (unit~102).}

\begin{proof}[Proof of Theorem~C]
Since $\sigma$ is a bijection, \m{|n'|=|n|+1}; write $N=|n|$. By the two hypotheses on the signs, all of
$\mathbf{s}$, $\mathbf{r}$, $\mathbf{t}$ are~$1$. Since $-l$ is a doubled Harish-Chandra parameter in rank~$N$, the
values of each $l_\tau$ are distinct and congruent to $N-1$ modulo~$2$; by~(I9), $l$ is the doubled Harish-Chandra
parameter of a weight $(\knat,\enat)$. For $0\le p\le N$ and $0\le q<N$ put
\[\lambda_\tau(p)=e_\tau+\sum_{p<h<N+1}\kappa_\tau(h),\qquad
\lambda^{\natural}_\tau(q)=\enat_\tau+\sum_{q<h<N}\knat_\tau(h),\]
the highest weights. The values of $l'_\tau$ in decreasing order are \m{2\lambda_\tau(p)+N-2p}, $p=0,\dots,N$, and
those of $l_\tau$ are \m{2\lambda^{\natural}_\tau(q)+N-1-2q}, $q=0,\dots,N-1$.

\emph{Interlacing.} Fix $\tau$ and $q<N$, and count values from~$0$. Exactly $q$ values of $l'_\tau$ exceed its
$q$-th value~$y$; by the hypothesis at this~$y$, at most $q$ values of $l_\tau$ exceed it, so the $q$-th value of
$l_\tau$ is at most the $q$-th value of $l'_\tau$. Exactly $q$ values of $l_\tau$ exceed its $q$-th value~$y$; by
the hypothesis at this~$y$, at most $q+1$ values of $l'_\tau$ exceed it, so the $(q+1)$-st value of $l'_\tau$ is
at most the $q$-th value of~$l_\tau$. The values of $l_\tau$ and of $l'_\tau$ have opposite parities, so both
inequalities are strict, and in terms of the highest weights they read
\begin{equation}\label{e:interlace}
\lambda_\tau(q+1)\ \le\ \lambda^{\natural}_\tau(q)\ \le\ \lambda_\tau(q)\qquad(0\le q<N).
\end{equation}
Put \m{\delta_\tau(q)=\lambda^{\natural}_\tau(q)-\lambda_\tau(q+1)}; it is half of the difference of the $q$-th value of
$l_\tau$ and the $(q+1)$-st value of $l'_\tau$, rounded down, and \m{0\le\delta_\tau(q)\le\kappa_\tau(q+1)}
by~\eqref{e:interlace}.

\emph{The matrix.} Let $\kdn_\tau(h)=\delta_\tau(h-1)$ and $\kup_\tau(h)=\kappa_\tau(h)-\delta_\tau(h-1)$ for
\m{1\le h\le N}, and $\kdn_\tau(h)=0$, $\kup_\tau(h)=\kappa_\tau(h)$ for the other~$h$. Then
\m{\kup+\kdn=\kappa}, and
\begin{align*}
\kdn_\tau(h)+\kup_\tau(h+1)&=\lambda^{\natural}_\tau(h-1)-\lambda^{\natural}_\tau(h)=\knat_\tau(h)\qquad(0<h<N),\\
e_\tau+\kdn_\tau(N)&=\lambda_\tau(N)+\delta_\tau(N-1)=\lambda^{\natural}_\tau(N-1)=\enat_\tau .
\end{align*}
So
(I11) gives a matrix~$\brm$, and (I10) a matrix~$A$. Put \m{B=A\cdot\brm} and
\[N_\tau=e_\tau+\sum_{0<h<N+1}\kup_\tau(h).\]

(a1) For $g\in\Un_n(S)$ and $u\in\Un_1(S)$, the intertwining property of $\brm$ over $R$, at the families $f(g)$
and $(f_\tau(\det u))_\tau$, and then the invariance of~$A$ give
\[\bigl\langle\wrep{\kappa'}{e'}(f(g))v',A\,\brm\,\wrep{\kappa}{e}(f(\dg{g,u}))v\bigr\rangle
=\prod_\tau f_\tau(\det u)^{N_\tau}\cdot\langle v',A\,\brm\,v\rangle .\]
(a2) \m{|N_\tau|\le|e_\tau|+\sum_{0<h<N+1}\kappa_\tau(h)\le|e_\tau|+\boxes{\kappa}{N+1}}.

(a3) By (I10), applied to the vector $\brm\,v$, there is $v^{\circ}$ with \m{\wn{v^{\circ}}=\wn{\brm\,v}} and
\m{\langle v',A\,\brm\,v\rangle=(v^{\circ},v')}. The Cauchy--Schwarz inequality for the positive definite Gram form and
(I11) give \m{|\langle v',Bv\rangle|\le\wn{\brm\,v}\,\wn{v'}\le(M^{M})^{3}\wn{v}\,\wn{v'}}, provided
\m{M\ge\boxes{\knat}{N}}. This holds for $M\ge\boxes{\kappa}{N+1}$: counting boxes by rows,
\[\boxes{\knat}{N}=\sum_\tau\sum_{q<N}(\lambda^{\natural}_\tau(q)-\enat_\tau),\qquad
\boxes{\kappa}{N+1}=\sum_\tau\sum_{p<N+1}(\lambda_\tau(p)-e_\tau),\]
and \m{\lambda^{\natural}_\tau(q)-\enat_\tau\le\lambda_\tau(q)-e_\tau} by~\eqref{e:interlace} and
\m{\enat_\tau\ge e_\tau}.

(b) Fix~$g$. Apply (I12) to the standard space, the matrix~$\brm$, the vector \m{v^{\natural}=v^{+}_{l}} and the
functional \m{L'=(v^{+}_{l'},\cj{\wrep{\kappa}{e}}(g)^{-1}\,\cdot\,)}; it gives~$C$. For $V$, $\eta$, $\phi$, $u$,
$\mathcal{K}$ as in the statement, $V$ is the standard space, $\Phi=\Fof{-l}{\mathbf{s}}{\eta}$ is automorphic, and by
(I10) it has the form required in (I12) with \m{L_h=\langle\eta(h),A\,\cdot\,\rangle}. The left side of (I12)
is then the integral of the theorem, because the value of \m{\Fof{l'}{\mathbf{t}}{\phi}} at \m{\dg{y,(1,u)}\cdot(g,1)} is
\m{L'\bigl(\cj{\wrep{\kappa}{e}}(\dg{y_\infty,1})^{-1}\phi(\dg{h,u})\bigr)}, and the right side is the mean of
\m{h\mapsto\langle\eta(h),A\,\brm\,\phi(\dg{h,u})\rangle}.
\end{proof}

\sect{The grid spaces}{s:grid}
Chapter~VII compares the chain of lifts with a second chain by moving through two grids of hermitian spaces, the
grid $V$ and the grid~$Z$. The file read here only names the objects on the grids; their index sets, signs and
parameters are defined in Chapter~VII. In this section $D$ is admissible data as in Setting~\ref{set:data}, and
$z\in\{0,1\}$ names a grid: $z=0$ is the grid~$V$, $z=1$ the grid~$Z$. (The formal text uses truth values; $1$ is
\emph{true}.)

\statement{ana}{\textbf{Inputs from Chapter VII} (units 140 and~147) \textbf{and from the background
library.}\enspace In Setting~\ref{set:data}:

(I13) \emph{Spaces.} A hermitian space with a point over $\ED{D}$, with a finite index set and signs on it, is a
Gram matrix over $\ED{D}$ with an adelic frame that brings it to the diagonal matrix of the signs at the infinite
places and to $1$ at the finite places. Such spaces have orthogonal sums $\oplus$, relabellings along bijections of
index sets that carry the signs, and opposites (the Gram matrix and the signs are negated); for signs $\mathbf{s}$
there is the standard signed space $I(\mathbf{s})$. For every natural number~$j$, $\SpY{D}{j}$ is a hermitian space
with index set $\Yidx{D}{j}$ and signs $\gsigns{0}{0}{0}{j}$ (unit~140).

(I14) \emph{Index sets, signs, parameters} (unit~147). For $z\in\{0,1\}$ and natural numbers $o,i,j$:
\m{\GrIdx{z}{i}{j}=\BlockIdx{j}{\BlockIdx{i}{\BlockIdx{z}{\Ridx{D}}{\{\pt\}}}{[\tD{D}+z]}}{[\tD{D}-z]}}, a set with
\m{\rD{D}+z+i\,(\tD{D}+z)+j\,(\tD{D}-z)} elements; its base indices, those of $\BlockIdx{z}{\Ridx{D}}{\{\pt\}}$, are
numbered by $\basisNum{D}$ and, for the point, by~$\rD{D}$. The signs $\gsigns{z}{o}{i}{j}$ on it are
\m{\esign{o}{j}{\max(q+1-\kD{D},0)}} at the base index of number~$q$ and $+1$ elsewhere. The doubled parameter
$\gpar{z}{o}{i}{j}{\tau}$, for an infinite place~$\tau$, is the integer vector on the same index set with the
following entries: $2\,\rowPar{D}{\tau}{q}$ at the base index of number~$q$;
\m{2\,\wt{D}{l+1}{\tau}+\tD{D}-1-2a} at the position $a\in[\tD{D}+z]$ of the block $l<i$ of the first kind;
\m{\tD{D}-1-2\,(\wt{D}{l+1+o}{\tau}+z)-2a} at the position $a\in[\tD{D}-z]$ of the block $l<j$ of the second kind;
all raised by $1$ if $z=1$, and lowered by $1$ if $z=0$ and $i+j$ is odd. The weight
$\gwt{c}{l}{n}$, for $c\in\{0,1\}$, a natural number $l$ and a finite set~$n$, is
\m{\tau\mapsto\lfloor(|n|+1-c)/2\rfloor-\wt{D}{l+1}{\tau}}. The polynomial \m{\ghar{o}{\delta}{c}{l}{z}{u}{v}{i}{j}}, for
$\delta\in\{0,1\}$, is a chosen joint harmonic for the nodes $(u,v)$ and $(i,j)$ of grid~$z$, the weights
$\gwt{c}{l}{\cdot}$ and the parameters, negated if $\delta=1$. Facts: (i) $\gsigns{1}{0}{0}{0}$ is $\gsigns{0}{0}{0}{0}$ on
$\Ridx{D}$ and $+1$ at the point; (ii) the number of signs $-1$ in $\gsigns{z}{0}{i}{j}$ is
\m{\lfloor(\max(\mD{D}-j,0)+z)/2\rfloor}; (iii) there are chosen bijections
\m{\Delta_j\colon\GrIdx{0}{j}{j}\sqcup\{\pt\}\to\GrIdx{1}{j}{j}}, the identity for $j=0$, such that for $j\le \mD{D}$
the signs $\gsigns{1}{0}{j}{j}$ correspond to $\gsigns{0}{0}{j}{j}$ and $(-1)^{j}$ at the point; (iv) for all
$z,i,j$ and~$\tau$, the entries of $\gpar{z}{0}{i}{j}{\tau}$ are distinct, congruent to $|\GrIdx{z}{i}{j}|-1$
modulo~$2$, and two entries at indices of different signs differ by at least $2\,\hD{D}$; (v) for a step of the
grid as in Lemma~\ref{l:edge}, the weights, parameters and signs at its two ends correspond in the sense of the theta
correspondence of discrete series parameters; (vi) $\tD{D}$ is odd.}

\statement{ana}{\textbf{Inputs from Chapter VII and from the background library, continued.}\enspace In
Setting~\ref{set:data}:

(I15) \emph{Moduli} (units 140 and~147). For a tame character $\vartheta$ of $D$ and a natural number $q\ge1$,
$\lmod{\vartheta}{q}$ is the ideal of $\OO_{\ED{D}}$ that is the $q$-th power of the product of the chosen places
of~$D$, times the product of the places of~$Q(\vartheta)$; its prime divisors are the chosen places and the places
of~$Q(\vartheta)$. For an $\mD{D}$-fold twist~$\vartheta$: \m{\gmod{\vartheta}=\lmod{\vartheta}{\bone{D}}},
\m{\tmod{0}{\vartheta}=\lmod{\vartheta}{1}} and \m{\tmod{1}{\vartheta}=\lmod{\oneT}{\bone{D}}}, which divides
$\gmod{\vartheta}$. For a finite index set $n$ and each of these ideals $\mathfrak{a}$, reduction from the stabiliser $\Kmax$ of the
standard lattice in $\Un_n(\Af)$ to \m{\GL_n(\OO_{\ED{D}}/\mathfrak{a})} is surjective.

(I16) \emph{Kernels and lifts} (background library). For a parameter $l$ and signs $\mathbf{s}$ on a finite index
set, the \kw{kernel pair} of $(l,\mathbf{s})$ sends $0$ to the product over the infinite places of the discrete series
kernels of the parameters~$l_\tau$, and $1$ to that of~$-l$. A pair of functions on the real unitary group is a
\kw{conjugate kernel pair} if the two are complex conjugates of each other and each is a cusp kernel of order~$0$.
If, for every~$\tau$, the entries of $l_\tau$ are distinct, of the right parity, and differ by at least
$1000\,|n|^3$ at indices of different signs, then the kernel pair of $(l,\mathbf{s})$ is a conjugate kernel pair.
For spaces $\mathcal{S}$, $\mathcal{T}$, characters of $\Un(1)$, a polynomial and a function $\Psi$ of matrices modulo an
ideal, there are the theta lift from forms on $\mathcal{S}$ to forms on $\mathcal{T}$ and the descent in the other
direction, both with the Schwartz function made of the polynomial at the infinite places and of $\Psi$ at the
finite places. A family of characters, two kernel pairs, a direction $\delta\in\{0,1\}$ and a polynomial form an
\kw{edge family} if both pairs are conjugate kernel pairs, the characters have weights independent of the member of
the family, and the polynomial is an edge for the real Weil representation between the kernels of index~$\delta$. If
weights, parameters and signs correspond as in (I14)(v), then the kernel pairs, the joint harmonic and any family of
characters of these weights form an edge family. \emph{Status.} The catalogue classes the two statements of (I16)
as substantial and gives no reference for them: reference to be confirmed. The second rests on the library's
statement that, for corresponding parameters, the vector of the Fock model given by the product of the chosen
joint harmonics is an edge between the two discrete-series types. For that statement the catalogue cites
Ichino~[I], Theorem~4.1 and Lemma~7.8, for the compact types, and records the verdict of its second reading: the
formal matching condition selects witnesses for a two-sided analytic statement beyond the published joint-type
criterion. So the formal statement is not a special case or a weaker form of the cited ones; it goes beyond them,
and its relation to the literature is unresolved. It is used here as the library proves it.

(I17) \emph{Determinants} (background library). If a character $\nu$ of $\Un(1)$ of $\ED{D}$ has a weight and a tame
part $\psi$ that is trivial outside the prime divisors of an ideal $\mathfrak{a}$, then for every finite index set $n'$,
all signs $\mathbf{s}$ on $n'$ and every $\mathbf{k}\in\Kmax\subseteq\Un_{n'}(\Af)$: the value of $\nu$ at the determinant of
the adelic point $(1,\mathbf{k})$ is \m{\prod_{w\mid\mathfrak{a}}\psi_w\bigl(\det(\mathbf{k}\bmod w)\bigr)}.}

\begin{definition}[nodes and their spaces]\label{d:grid}
In Setting~\ref{set:data}, let $D$ be admissible data.
\begin{enumerate}[label=(\alph*)]
\item A \kw{node} is a triple $p=(z,u,v)$ of an element $z\in\{0,1\}$ and natural numbers $u$, $v$.
\item For a node $p=(z,i,j)$: \m{\nodeIdx{p}=\GrIdx{z}{i}{j}}.
\item For natural numbers $o,i,j$ and $z\in\{0,1\}$, the hermitian space $\gridSp{o}{z}{i}{j}$ has index set
$\GrIdx{z}{i}{j}$ and signs $\gsigns{z}{o}{i}{j}$. It is: \m{\SpY{D}{0}\oplus I(+1)} for $(o,z,i,j)=(0,1,0,0)$, where
$I(+1)$ is the standard line of sign~$+1$; $\SpY{D}{j}$ for $(o,z,i)=(0,0,0)$; the opposite of
$\gridSp{o-1}{z}{0}{0}$ for $o\ge1$ and $i=j=0$; and the standard signed space $I(\gsigns{z}{o}{i}{j})$ in all other
cases.
\end{enumerate}
\end{definition}
\leanline{\lean{Node}\sep \lean{NodeIndex}\sep \lean{gridSpace}}
The first case has the right signs by (I14)(i), and the third because
$\esign{o}{0}{a}=-\esign{o-1}{0}{a}$; this is why Definition~\ref{d:esign} begins with that equation.

\begin{definition}[forms on the grids]\label{d:gridforms}
In Setting~\ref{set:data}, let $D$ be admissible data, $z\in\{0,1\}$ and $i,j,u,v$ natural numbers.
\begin{enumerate}[label=(\alph*)]
\item $\gridKer{z}{i}{j}$ is the kernel pair of the parameter \m{\tau\mapsto\gpar{z}{0}{i}{j}{\tau}} and the signs
$\gsigns{z}{0}{i}{j}$ (Input~(I16)).
\item $\gridForm{z}{u}{v}$ is the space of the forms on the unitary group of $\gridSp{0}{z}{u}{v}$ whose type is the
kernel of index $1$ of $\gridKer{z}{u}{v}$, that of the negated parameter; it carries the Petersson product.
\item $\gridLevForms{i}$ is the subspace of $\gridForm{1}{0}{i}$ of the forms fixed by the principal level of the
integer \m{\piGbar{D}^{\,\bone{D}}}.
\end{enumerate}
\end{definition}
\leanline{\lean{gridKernel}\sep \lean{GridForm}\sep \lean{gridLevelForms}}

\begin{definition}[characters on the grids]\label{d:gridchar}
In Setting~\ref{set:data}, let $D$ be admissible data, $\vartheta$ an $\mD{D}$-fold twist of~$D$, $c\in\{0,1\}$, $l$ a
natural number and $n$ a finite index set.
\begin{enumerate}[label=(\alph*)]
\item the character \m{\gridChar{\vartheta}{c}{l}{n}} is given by
\[\gridChar{\vartheta}{c}{l}{n}=\begin{cases}
\liftChar{D}{\vartheta^{l+1}}{l+1}\cdot\ups^{\lfloor|n|/2\rfloor}&\text{if }c=1,\\[2pt]
\liftChar{D}{\oneT}{l+1}\cdot\ups^{\lfloor(|n|+1)/2\rfloor}&\text{if }c=0.\end{cases}\]
\item \m{\gridTame{\vartheta}{c}{l}=\liftTame{D}{\vartheta^{l+1}}{l+1}} if $c=1$, and
\m{\gridTame{\vartheta}{c}{l}=\liftTame{D}{\oneT}{l+1}} if $c=0$.
\end{enumerate}
\end{definition}
\leanline{\lean{gridChar}\sep \lean{gridTamePart}}

\begin{definition}[levels on the grids]\label{d:gridlevels}
In Setting~\ref{set:data}, let $D$ be admissible data, $\vartheta$ an $\mD{D}$-fold twist of~$D$, $z\in\{0,1\}$ and $u$,
$v$ natural numbers.
\begin{enumerate}[label=(\alph*)]
\item $\triOrd{z}{\vartheta}{v}$ is the order of the group
\m{\GL_{\GrIdx{z}{0}{v}}\bigl(\OO_{\ED{D}}/\tmod{z}{\vartheta}\bigr)}.
\item For a node $p=(z,u,v)$, $\nodeLev{\vartheta}{p}$ is the following subgroup of the group
\m{\GL_{\nodeIdx{p}}\bigl(\OO_{\ED{D}}/\gmod{\vartheta}\bigr)}: it is trivial if $z=0$ and $u=0$, and otherwise it is
the kernel of reduction to \m{\GL_{\nodeIdx{p}}\bigl(\OO_{\ED{D}}/\tmod{1}{\vartheta}\bigr)}.
\item $\levIdx{\vartheta}{z}{u}{v}$ is the index of $\nodeLev{\vartheta}{(z,u,v)}$ in that group, as a complex
number.
\end{enumerate}
\end{definition}
\leanline{\lean{triangleOrder}\sep \lean{nodeLevel}\sep \lean{levelIndex}}

\begin{definition}[moves on the grids]\label{d:moves}
In Setting~\ref{set:data}, let $D$ be admissible data, $\vartheta$ an $\mD{D}$-fold twist of~$D$, $z\in\{0,1\}$ and $u$,
$v$ natural numbers. In each case $F$ is a form in the source of the map.
\begin{enumerate}[label=(\alph*)]
\item \m{\ptDesc{\vartheta}{z}{v}{\,\cdot\,}\colon\gridForm{z}{0}{v+1}\to\gridForm{z}{0}{v}} is
$\triOrd{z}{\vartheta}{v}$ times the descent (Input~(I16)) from the space $\gridSp{0}{z}{0}{v+1}$ to the space
$\gridSp{0}{z}{0}{v}$, formed with the characters \m{n\mapsto\gridChar{\vartheta}{1-z}{v}{n}}, with the polynomial
\[\ghar{0}{0}{1-z}{v}{z}{0}{v}{0}{v+1}\]
and, at the finite places, with the indicator of the matrix of the inclusion of $\GrIdx{z}{0}{v}$ into
$\GrIdx{z}{0}{v+1}$ modulo~$\tmod{z}{\vartheta}$; we write $\ptDesc{\vartheta}{z}{v}{F}$ for its value at~$F$.
\item \m{\iotaPlus{\vartheta}{z}{v}{\,\cdot\,}\colon\gridForm{z}{0}{\mD{D}}\to\gridForm{z}{0}{v}} is defined by descending
recursion: it is the identity for $v=\mD{D}$, zero for $v>\mD{D}$, and for $v<\mD{D}$
\[\iotaPlus{\vartheta}{z}{v}{F}=\ptDesc{\vartheta}{z}{v}{\iotaPlus{\vartheta}{z}{v+1}{F}}.\]
\item For a matrix $y$ over \m{\OO_{\ED{D}}/\gmod{\vartheta}} with rows indexed by $\GrIdx{z}{u}{v+1}$ and columns by
$\GrIdx{z}{u}{v}$, $\descMove{y}{\,\cdot\,}\colon\gridForm{z}{u}{v+1}\to\gridForm{z}{u}{v}$ is the descent with the
characters $n\mapsto\gridChar{\vartheta}{1-z}{v}{n}$, the polynomial $\ghar{0}{0}{1-z}{v}{z}{u}{v}{u}{v+1}$ and
the indicator of $y$ modulo~$\gmod{\vartheta}$.
\item For a matrix $y$ over \m{\OO_{\ED{D}}/\gmod{\vartheta}} with rows indexed by $\GrIdx{z}{u+1}{v}$ and columns by
$\GrIdx{z}{u}{v}$, $\liftMove{y}{\,\cdot\,}\colon\gridForm{z}{u}{v}\to\gridForm{z}{u+1}{v}$ is the lift with the
characters $n\mapsto\gridChar{\vartheta}{z}{u}{n}$, the polynomial $\ghar{0}{1}{z}{u}{z}{u}{v}{u+1}{v}$ and the
indicator of $y$ modulo~$\gmod{\vartheta}$.
\end{enumerate}
\end{definition}
\leanline{\lean{pointDescent}\sep \lean{iotaPlus}\sep \lean{descentMove}\sep \lean{liftMove}}

\begin{lemma}\label{l:gridkernel}
In Setting~\ref{set:data}, let $D$ be admissible data, $z\in\{0,1\}$ and $i$, $j$ natural numbers with $i\le \mD{D}$ and
$j\le \mD{D}$. Then $\gridKer{z}{i}{j}$ is a conjugate kernel pair.
\end{lemma}
\begin{proof}
The index set has at most \m{2\,\nD{D}{\mD{D}}} elements (this count is in unit~147), and
\[\nD{D}{\mD{D}}=\kD{D}+\mD{D}\,(5\rD{D}+1)\le \rD{D}+12\sqrt{\xD{D}}\,\rD{D}\le2\rD{D}^2\]
by~(D2). By (I14)(iv) the entries of the parameter at each $\tau$ are distinct, of the right parity, and differ by at
least \m{2\,\hD{D}=4\rD{D}^8\ge1000\,(4\rD{D}^2)^3} at indices of different signs; so (I16) applies.
\end{proof}
\leanline{\lean{isConjKernelPairgridKernel}}

\begin{lemma}\label{l:gridcharweight}
In Setting~\ref{set:data}, let $D$ be admissible data, $\vartheta$ an $\mD{D}$-fold twist of~$D$, $c\in\{0,1\}$, $l$ a
natural number and $n$ a finite index set.
\begin{enumerate}[label=(\roman*)]
\item $\gridChar{\vartheta}{c}{l}{n}$ has weight $\gwt{c}{l}{n}$ and tame part $\gridTame{\vartheta}{c}{l}$.
\item Let $n'$ be a finite index set, $\mathbf{s}$ signs on $n'$ and $\mathbf{k}$ an element of the stabiliser $\Kmax$ of the
standard lattice in $\Un_{n'}(\Af)$. Then the value of $\gridChar{\vartheta}{c}{l}{n}$ at the determinant of the
adelic point $(1,\mathbf{k})$ of the unitary group of the signs $\mathbf{s}$ is
\[\prod_{w\mid\gmod{\vartheta}}\gridTame{\vartheta}{c}{l}_w\bigl(\det(\mathbf{k}\bmod w)\bigr),\]
the value at \m{\det(\mathbf{k}\bmod\gmod{\vartheta})} of the character of \m{(\OO_{\ED{D}}/\gmod{\vartheta})^\times} with these
components. In particular it does not depend on~$n$.
\end{enumerate}
\end{lemma}
\begin{proof}
(i) By Theorem~B(iii) the first factor has weight \m{-\wt{D}{l+1}{\tau}} and the stated tame part; by (I8) and (I7),
multiplying by $\ups^{y}$ adds $y$ to the weight. (ii) By Lemma~\ref{l:omega}(ii) and Definition~\ref{d:tame}, the
tame part is trivial outside the chosen places and the places of $Q(\vartheta)$, which by (I15) are the prime
divisors of $\gmod{\vartheta}$; apply (I17) with~(i).
\end{proof}
\leanline{\lean{gridCharhasWeight}\sep \lean{gridChardet}}

\begin{lemma}[steps of the grids are edges]\label{l:edge}
In Setting~\ref{set:data}, let $D$ be admissible data, $z\in\{0,1\}$, $i,j$ natural numbers, $I$ a set,
\m{\vartheta=(\vartheta_y)_{y\in I}} a family of $\mD{D}$-fold twists of~$D$ and $\delta\in\{0,1\}$. Let $c\in\{0,1\}$ and
natural numbers $l,i',j'$ be given by \m{(c,l,i',j')=(z,i,i+1,j)} if $\delta=1$ and \m{(c,l,i',j')=(1-z,j,i,j+1)} if
$\delta=0$, and suppose $i'\le \mD{D}$ and $j'\le \mD{D}$. Then the characters
\m{n\mapsto\gridChar{\vartheta_y}{c}{l}{n}}, $y\in I$, the kernel pairs $\gridKer{z}{i}{j}$ and
$\gridKer{z}{i'}{j'}$, the direction $\delta$ and the polynomial \m{\ghar{0}{\delta}{c}{l}{z}{i}{j}{i'}{j'}} form an edge
family.
\end{lemma}
\begin{proof}
By Lemma~\ref{l:gridcharweight}(i) the characters have the weights $\gwt{c}{l}{\cdot}$, which do not depend on~$y$.
By (I14)(v) weights, parameters and signs at the two nodes correspond, and (I16) gives the edge family.
\end{proof}
\leanline{\lean{isEdgeFamilygrid}}

\begin{lemma}[the two grids meet]\label{l:diag}
In Setting~\ref{set:data}, let $D$ be admissible data and $j$ a natural number.
\begin{enumerate}[label=(\roman*)]
\item If $j\le \mD{D}$, then $\gridSp{0}{1}{j}{j}$ is the orthogonal sum of $\gridSp{0}{0}{j}{j}$ and the standard line of
sign $(-1)^j$, relabelled along the bijection $\Delta_j$ of~(I14)(iii).
\item There is a bijection from $\GrIdx{0}{j+1}{j}$ to $\GrIdx{1}{j}{j+1}$ that carries the signs
$\gsigns{0}{0}{j+1}{j}$ to $\gsigns{1}{0}{j}{j+1}$ and along which $\gridSp{0}{1}{j}{j+1}$ is the relabelling of
$\gridSp{0}{0}{j+1}{j}$.
\end{enumerate}
\end{lemma}
\begin{proof}
(i) For $j=0$ both sides are \m{\SpY{D}{0}\oplus I(+1)} by Definition~\ref{d:grid}, and $\Delta_0$ is the
identity. For $j\ge1$ both $\gridSp{0}{1}{j}{j}$ and $\gridSp{0}{0}{j}{j}$ are standard signed spaces. A sum of
standard signed spaces is the standard signed space of the two families of signs, and its relabelling is the
standard signed space of the transported signs, which are $\gsigns{1}{0}{j}{j}$ by~(I14)(iii).

(ii) Both index sets have \m{\rD{D}+(2j+1)\,\tD{D}} elements by (I14), since $\tD{D}\ge1$. By (I14)(ii) the numbers of
signs $-1$ are \m{\lfloor(\max(\mD{D}-j-1,0)+1)/2\rfloor} and \m{\lfloor\max(\mD{D}-j,0)/2\rfloor}, which agree. So there is a
bijection carrying one family of signs to the other. Both spaces are standard signed spaces by
Definition~\ref{d:grid}, since $j+1\ne0$; relabelling one gives the other.
\end{proof}
\leanline{\lean{gridSpacediag}\sep \lean{gridSpacesubdiag}}

\begin{lemma}\label{l:levelindex}
In Setting~\ref{set:data}, let $D$ be admissible data, $\vartheta$ an $\mD{D}$-fold twist of~$D$, $z\in\{0,1\}$ and $u,v$
natural numbers, and suppose that $z=1$ or $u\ne0$. Then
\[\levIdx{\vartheta}{z}{u}{v}=\bigl|\GL_{\GrIdx{z}{u}{v}}\bigl(\OO_{\ED{D}}/\gmod{\oneT}\bigr)\bigr| .\]
In particular it does not depend on~$\vartheta$.
\end{lemma}
\begin{proof}
Under the hypothesis, $\nodeLev{\vartheta}{(z,u,v)}$ is the kernel of the reduction
\[\GL\bigl(\OO_{\ED{D}}/\gmod{\vartheta}\bigr)\longrightarrow\GL\bigl(\OO_{\ED{D}}/\tmod{1}{\vartheta}\bigr),\]
and $\tmod{1}{\vartheta}=\gmod{\oneT}$ by~(I15). The index of a kernel is the order of the image, and the reduction
is surjective: by (I15) reduction from $\Kmax$ is surjective and factors through it.
\end{proof}
\leanline{\lean{levelIndexequntwisted}}

\sect{Where the results are used}{s:use}
The objects of Sections \ref{s:chosen} to \ref{s:index} are used in every later chapter. The level $\levelN{D}$ and
the element $\piN{D}$ define the level of the columns of Chapter~V (unit~145) and the split places where the
Gauss-sum operator of the note \emph{Chains of definite theta lifts and the Gauss-sum operator} acts; the lower bound
\m{\log\levelN{D}\ge\logTarget{D}} of Theorem~A(ii) is what lets a power of $\levelN{D}$ pay for the excess of size
over divisibility in the last step of the theorem on short forms (unit~151), and the form of $\logTarget{D}$
enters the admissibility of the filling data (unit~148). The characters $\liftChar{D}{\vartheta}{j}$ are those of the $j$-th
theta lift in the seesaw identity and in Chapter~VII; the index sets $\Xidx{D}{j}$ and $\Yidx{D}{j}$ carry the
definite and the indefinite chain. The group $\Ggrp{D}$ and the idempotent $\ezero{D}$ define the rows (unit~140).
Pairs are the arguments of the ratios of $L$-values of Chapters VII and~VIII. Theorem~C and everything of
Section~\ref{s:grid} are used only in the folder of the exchange identity (unit~147); so are entry signs and
Lemma~\ref{l:ite}.

\sect{Notes on the formal text}{s:notes}
\subhead{Arguments that the census lists and Lean does not take.} The lists of arguments used by the build of this
note are read from the text of the source, not printed by Lean. In the file of index types a line of the form
\emph{variable} declares a finite place~$w$ and natural numbers $a$, $n$, $j$ for the whole file, and the reading
attributes such a variable to a declaration whenever its name occurs in it, also as the name of a field or of a bound
variable. So it lists $w$ among the arguments of the definitions of $\ChPl{D}{i}$, $\PairSet{D}$, the row parameter,
$\TameSet{D}$, $\levelNorm{\vartheta}$, $\TwistSet{D}$, $\goodFam{D}$ and of statements of Theorem~B and of
Lemmas \ref{l:chosen}(v) and~\ref{l:muflat}; a type $T$ in Lemma~\ref{l:reschar} (it is the character there); a number $a$
in Lemma~\ref{l:sizes}(i); a number $n$ in the definition of $\InitIdx{D}{j}$ (the name occurs as the name of the
function $\nD{D}{j}$); two numbers $a$, $j$ in the definition of entry signs (they are names of pattern variables);
an index set $n$ in the definition of $\nodeLev{\vartheta}{p}$ (it names an argument of another function); and
the field with its structures and a finiteness in Lemma~\ref{l:gridcharweight}(ii). Lean takes none of these. For
the last three definitions this note therefore writes the notation with Lean's arguments ($D$ and $j$; $o$, $j$,
$a$; $\vartheta$ and $p$) and not with the census's. In the other direction, the reading takes the source of each
of the four linear maps of Definition~\ref{d:moves} for an argument; the notation shows the form to which the map is
applied in that place. The arguments of $\BlockIdx{j}{S}{T}$ are, in Lean's order, $S$, $T$ and the number~$j$.

\subhead{Implicit data in Section~\ref{s:grid}.} In the file of the grid spaces the data $D$ are an implicit
argument of every definition, so the notations of that section show $D$ as a fixed letter and not in a slot of
their own; they are to be read with the $D$ of Setting~\ref{set:data}.

\subhead{Where the formal statements differ from the ones above.}
Theorem~A(iii): the formal proof takes the square root of $-1$ from a lemma of the background library on fields
containing a root of unity of order $5040$; the proof above takes it from the root of unity of~$\ED{D}$.
Lemma~\ref{l:chosen}(iii) repeats ``$w$ is chosen'' on its right-hand side, as the formal statement does; its content
is that a chosen place lies above $g_0$ or~$g_1$. Lemma~\ref{l:rowindex}: the formal lemma is the count
\m{|\Ridx{D}|=\rD{D}}; the other two statements are added here. Definition~\ref{d:rowpar}(b): the formal text divides
$\kD{D}$ by $2$ as an integer, which is exact; the index $i+2-\kD{D}$ is a truncated difference of natural numbers,
used only for $i+1\ge\kD{D}$. Definition~\ref{d:pair}: the formal structure names its components
\emph{character}, $c$, $\psi$, $z$; the bound $4\bD{D}$ is a real number. Definition~\ref{d:tame}(a): the coprimality is
stated for the natural numbers \m{\gD{D}\cdot\levelN{D}} and \m{\prod_{w\in Q}\Nm w}. Theorem~C: the hypotheses on
the signs force all three families to be~$1$, so the theorem concerns definite groups only; the letter for the sign
on the point is an argument of the formal statement although it is determined. In (a1) the formal statement
quantifies over all rings $S$ with an involution and all fields $R$ of characteristic~$0$, as above; the
expression \m{\langle v',Bv\rangle} is bilinear there, in (a3) and in the mean of~(b). In (b) the constant $C$
depends on $g$ only: it is chosen after $g$ and before the space, the forms, $u$ and the level; the statement does
not say that $C\ne0$. Input~(I11) is
the case $q=0$ of a formal statement with one more parameter. Lemma~\ref{l:edge}: the two bounds $i'\le \mD{D}$,
$j'\le \mD{D}$ are arguments with a default proof by arithmetic. Lemma~\ref{l:diag}(i): the hypothesis $j\le \mD{D}$ enters only
through the signs of~(I14)(iii).

\subhead{Members that nothing refers to by name.} The reader's dossier lists a few members of the unit to which
no other declaration refers by name. Those of this note are: the statement $\pi_g\ne0$ of
Lemma~\ref{l:chosen}(i), which Lean's simplifier uses, and the bound for the root discriminant
(Lemma~\ref{l:sizes}(ii)) and the finiteness of $\BlockIdx{j}{S}{T}$ (Lemma~\ref{l:block}), which are instances found
by Lean's search. The instances of the file of index types that have no name have no line in the concordance;
they are mentioned in Lemmas \ref{l:sizes}(iii) and~\ref{l:chosen}(v), after Lemma~\ref{l:rowindex}, in
Remark~\ref{r:levelsizes} and after Definition~\ref{d:tame}. Whether a hypothesis of a formal statement is unused was not determined: no Lean
was run.

\subhead{How this note was written.} From the dossier of the unit: the texts of its members and the
structured proofs of the declarations treated here. I read the two Lean files of this note completely, the proof
of Theorem~C line by line, and compared every statement of this note with the Lean statement, argument by
argument, with the lists of arguments of the census and of the dossier. For the Inputs I read the Lean statements
of (I6), (I9) to (I12) and the definitions behind (I13) to (I17) in the files of Chapters III and~VII and of the
background library; I opened none of their proofs. The proofs of Lemma~\ref{l:sizes}, of Theorem~A(i) and of
Lemma~\ref{l:gridkernel} replace Lean's automatic comparisons of sizes by explicit inequalities from (D1) to~(D3),
and are mine; so are the proof of Theorem~A(iii) and the wording of the interlacing in the proof of Theorem~C, which
the formal text does by a case analysis left to automation. Two step claims of the dossier were corrected against
the Lean text: in Theorem~B(i) the residue characters are trivial at the root of unity (the dossier says that their
order divides $\Mom{P}{\mu}$), and in Theorem~C(a3) the vector of equal norm is attached to $\brm\,v$, not to~$v'$.
The program check of the unit's dossier covered \UnitSteps{} leaf steps of \UnitProofs{} proofs.

\subhead{Revision.} After an independent referee's report and the second reading of the catalogue, Inputs (I6)
and (I16) state the status of the cited literature as the catalogue records it (in both cases the relation is
unresolved), and the index of the middle vector of the shape is written $\diamond$, apart from the point~$\pt$.

\head{Appendix A. Concordance with the Lean artifact}
The formal names of the objects: \texttt{P.N₀}, \texttt{P.d} are $N_0$, $d$; the fields \texttt{D.logX},
\texttt{D.logB}, \texttt{D.s}, \texttt{D.m}, \texttt{D.r}, \texttt{D.g}, \texttt{D.piG} are $\log \xD{D}$, $\log \bD{D}$,
$\mu$, $\mD{D}$, $\rD{D}$, $(g_0,g_1)$, $\pi_g$; \texttt{D.Φ} is $\PlD{D}$, \texttt{D.w j τ} is $\wt{D}{j}{\tau}$,
\texttt{D.nominalWeight} the nominal weight, \texttt{D.gProd} is $\gD{D}$, \texttt{P.Mω D.s} is $\Mom{P}{\mu}$;
\texttt{star} is complex conjugation; \texttt{w.q} is $\Nm w$; \texttt{AutChar} is the group of characters of
$\Un(1)$ and \texttt{HasWeight} the relation ``has weight and tame part''; \texttt{WeilData.upsilonU1} is~$\ups$;
truth values are written $1$ and~$0$. Three statements as they stand in the artifact:
\quoted{TameChar}
\quoted{decld}
\quoted{gridSpace}

{\fontsize{8.6}{11.6}\selectfont\renewcommand{\leanfont}{\ttfamily\fontsize{7.7}{10}\selectfont\color{leancol}}
\setlength{\LTleft}{0pt}\setlength{\LTright}{\fill}\setlength{\tabcolsep}{0pt}
\begin{longtable}{@{}p{30mm}p{56mm}p{43mm}r@{}}
\textbf{In this note} & \textbf{Lean declaration} & \textbf{File} & \textbf{Lines}\\\midrule\endhead
\ConcordanceRows
\end{longtable}}

\head{Notation}
\begingroup\fontsize{8.6}{10.9}\selectfont\raggedright
\setlength{\columnsep}{9mm}\setlength{\parskip}{0pt}
\begin{multicols}{2}
\nt{$P=(N_0,d)$, $\KP{P}$}{parameters; the field \m{\QQ(\sqrt{-d})}}
\nt{$d$}{the prime of the parameters $P=(N_0,d)$, $d\equiv N_0/2-1$ modulo $N_0$ (Convention~0.1)}
\nt{$D$}{admissible data for~$P$ (Setting~\ref{set:data})}
\nt{$\xD{D}$, $\xoD{D}$, $\bD{D}$}{the three scales of~$D$}
\nt{$\mu$, $\mD{D}$, $\rD{D}$}{the exponent, the number of lifts and the rank of~$D$}
\nt{$g_0$, $g_1$, $\gD{D}$, $\pi_g$}{the two primes of~$D$, their product, an integer of $\KP{P}$ of norm $g_0g_1$}
\nt{$\cP{P}$}{the constant of~$P$ of Chapter~III}
\nt{$\ED{D}$, $\None{D}$, $\rtu{D}$}{the CM field of~$D$; the order of its root of unity; that root}
\nt{$\PlD{D}$, $\tau$, $\fD{D}$}{the infinite places of $\ED{D}$; one of them; their number}
\nt{$\thick{D}$}{the thick places}
\nt{$w$, $\cj{w}$, $\FF_w$, $\Nm w$}{a finite place of $\ED{D}$; its conjugate; residue field; absolute norm}
\nt{$\OO$, $\Af$}{rings of integers; the finite adeles of the field in question}
\nt{$\rdD{D}$, $\vpi{P}{\mu}$, $\OmP{P}{\mu}$, $\chow{d}$}{root discriminant; generator of the different; the period;
a constant of~$d$}
\nt{$\Mom{P}{\mu}$}{the number \m{(N_0d)^{\mu+6}}}
\nt{$\kD{D}$, $\tD{D}$, $\hhD{D}$, $\hD{D}$, $\nD{D}{j}$}{the numbers $r+1-m$, $5r$, $r^8$, $2r^8$,
\m{\kD{D}+j(\tD{D}+1)}}
\nt{$\aD{D}$, $\bone{D}$, $\btwo{D}$}{least multiples of $\None{D}$ above $xb$, $\aD{D}/x$, $\aD{D}/x_1$}
\nt{$\wt{D}{j}{\tau}$, $\nw{D}{\tau}$}{the weights; the nominal weight}
\nt{$\piGbar{D}$}{the conjugate of $\pi_g$ in $\OO_{\ED{D}}$}
\nt{$\isChosen{D}{w}$, $\notChosen{D}{w}$}{$w$ is, is not, a chosen place of~$D$}
\nt{$\ChPl{D}{i}$, $\chosenOf{D}{x}$}{chosen places above $g_i$; the place of a pair $x=(i,w)$ (see $x$ below)}
\nt{$x$}{in Definition~\ref{d:chosen} and Definition~\ref{d:G} also a pair $(i,w)$ of an index $i\in\{0,1\}$ and a chosen place $w$ (the scale of $D$ is written $\xD{D}$)}
\nt{$\omD{D}{w}$, $\omD{D}{\bullet}$, $T_w$}{residue character at~$w$; the family; the set of characters of order $\Mom{P}{\mu}^2$}
\nt{$\Un(1)$, $\Un_n$}{the norm-one group; the unitary group of the identity matrix}
\nt{$\nu$, $\mathbf{e}$, $\psi$}{a character of $\Un(1)$ (in Definition~\ref{d:block}, a bijection); a weight; a
tame part}
\nt{$\ups$}{the splitting character on~$\Un(1)$}
\nt{$\muFlat{D}$, $\chiCirc{D}$, $\muZero{D}{j}$}{the three characters of Definition~\ref{d:muflat}}
\nt{$\PairSet{D}$, $(\nu,\mathbf{c},\psi,z)$}{the pairs of~$D$; a pair}
\nt{$\logTarget{D}$, $\levelN{D}$, $\piN{D}$}{the target; the level; an integer of norm $\levelN{D}$}
\nt{$U$, $p$, $M_0$, $\Psi$, $\Psi_0$, $\rtuz$}{in Section~\ref{s:level}: a real number, a prime, $\None{D}/N_0$,
cyclotomic polynomials of orders $\None{D}$ and~$N_0$, a root of unity of order~$N_0$; $p$ also a node of the grids in Section~\ref{s:grid}}
\nt{$q$, $y$, $y'$}{natural numbers and integers bound in a statement}
\nt{$\TameSet{D}$, $\TwistSet{D}$, $\oneT$}{tame characters and twists of~$D$; the trivial one}
\nt{$\vartheta$, $Q(\vartheta)$, $\vartheta^{\mathrm{t}}$, $\levelNorm{\vartheta}$}{a tame character; its set of places; its
tame part; the product of the norms of the places of $Q(\vartheta)$}
\nt{$\liftChar{D}{\vartheta}{j}$, $\liftTame{D}{\vartheta}{j}$, $\LiftIdx{D}$}{character and tame part of the
$j$-th lift; the set $\{1,\dots,m\}$}
\nt{$[q]$, $\{\pt\}$}{the set $\{0,\dots,q-1\}$; a set with one element}
\nt{$S$, $T$, $\alpha$, $\BlockIdx{j}{S}{T}$}{sets; $S$ and $j$ copies of~$T$; $S$ also a ring in Theorem~C and in Section~\ref{s:branching}; $T$ also a set of characters in Sections~\ref{s:chosen} and~\ref{s:chars}}
\nt{$\firstBlock{j}$, $\equivSum{\nu}{j}$}{inclusion of the first block; a bijection}
\nt{$\rhoD{D}$}{the number $m/2-1$}
\nt{$\InitIdx{D}{j}$, $\Xidx{D}{j}$, $\Ridx{D}$, $\Yidx{D}{j}$}{index sets: of $X_j(D)$ without its last vector, of
$X_j(D)$, of $\SpY{D}{0}$, of $\SpY{D}{j}$}
\nt{$\shapeD{D}$}{the shape of $\SpY{D}{0}$}
\nt{$\basisNum{D}$}{a numbering of $\Ridx{D}$}
\nt{$\rowPar{D}{\tau}{i}$, $\parW{D}{\tau}{i}$}{the row parameter; its double on $\Ridx{D}$}
\nt{$\signsW{D}{j}{i}$, $\esign{o}{j}{a}$}{signs of the spaces $W_j$; entry signs}
\nt{$\Ggrp{D}$, $\goodFam{D}$, $\ezero{D}$, $c_x$, $e_{c_x}$}{the finite group; good families; the idempotent;
a polynomial and its idempotent}
\nt{$E$, $n$, $n'$}{in Section~\ref{s:branching}: a CM field, index sets}
\nt{$N$}{in Section~\ref{s:branching} also the number $|n|$ (the level $\levelN{D}$ is written in full)}
\nt{$(\kappa,e)$, $(\kappa',e')$, $(\knat,\enat)$}{weights}
\nt{$\kup$, $\kdn$}{the two parts of $\kappa$ in (I11)}
\nt{$l$, $l'$, $\hc{\kappa}{e}{N}{\tau,p}$}{families of parameters; the doubled Harish-Chandra parameter; $l$ also a natural number in Section~\ref{s:grid}}
\nt{$\lambda_\tau(p)$, $\lambda^{\natural}_\tau(q)$, $\delta_\tau(q)$}{highest weights; their difference}
\nt{$\mathbf{s}$, $\mathbf{r}$, $\mathbf{t}$, $\sigma$}{signs; a bijection \m{n\sqcup\{\pt\}\to n'}}
\nt{$\wrep{\kappa}{e}$, $\cj{\wrep{\kappa}{e}}$}{the representation of weight $(\kappa,e)$; at conjugate matrices}
\nt{$\langle v',v\rangle$, $(v',v)$, $\wn{v}$}{dot product; Gram form; Weyl norm}
\nt{$\boxes{\kappa}{N}$}{the number of boxes of $\kappa$ in rank~$N$}
\nt{$R$, $f_\tau$, $g$, $h$}{rings; ring homomorphisms; group elements}
\nt{$B$, $A$, $\brm$, $N_\tau$, $M$, $C$}{integer matrices; exponents; a natural number; a constant}
\nt{$\eta$, $\phi$, $\Phi$, $F$}{algebraic forms; an automorphic function; a form}
\nt{$\Fof{l}{\mathbf{s}}{\eta}$, $v^{+}_{l}$, $v^{\natural}$, $v^{\circ}$}{the function of a form; highest vector;
vectors}
\nt{$\mathcal{K}$, $\Kmax$, $\clint{\mathcal{K}}{F}$}{a level; the maximal one; the mean at level $\mathcal{K}$}
\nt{$\dg{g,u}$}{block diagonal matrix along~$\sigma$}
\nt{$L'$, $L_h$}{linear functionals}
\nt{$z$, $c$, $\delta$}{elements of $\{0,1\}$: a grid, twisted or not, a direction; $z$ is a natural number in a pair,
$c$ a natural number in Theorem~B and Lemma~\ref{l:chosen}}
\nt{$o$, $i$, $j$, $a$, $u$, $v$, $i'$, $j'$}{natural numbers: an offset and positions on the grids, $u$ and $v$ the
offsets of a node; $u$ also an element or adele of norm one in Section~\ref{s:chars}}
\nt{$I(\mathbf{s})$, $\SpY{D}{j}$}{standard signed space; the hermitian space $Y_j(D)$}
\nt{$\GrIdx{z}{i}{j}$, $\gsigns{z}{o}{i}{j}$, $\gpar{z}{o}{i}{j}{\tau}$}{index set, signs and doubled parameter of a
node}
\nt{$\gwt{c}{l}{n}$, $\ghar{o}{\delta}{c}{l}{z}{u}{v}{i}{j}$}{grid weight; joint harmonic}
\nt{$\lmod{\vartheta}{q}$, $\gmod{\vartheta}$, $\tmod{z}{\vartheta}$, $\mathfrak{a}$}{moduli; an ideal}
\nt{$\nodeIdx{p}$}{the index set of a node $p=(z,u,v)$ (Section~\ref{s:grid})}
\nt{$\gridSp{o}{z}{i}{j}$}{the space at a node}
\nt{$\gridKer{z}{i}{j}$, $\gridForm{z}{u}{v}$, $\gridLevForms{i}$}{kernel pair; forms at a node; forms of level}
\nt{$\gridChar{\vartheta}{c}{l}{n}$, $\gridTame{\vartheta}{c}{l}$}{character and tame part on the grids}
\nt{$\triOrd{z}{\vartheta}{v}$, $\nodeLev{\vartheta}{p}$, $\levIdx{\vartheta}{z}{u}{v}$}{an order; the level at a
node; its index}
\nt{$\ptDesc{\vartheta}{z}{v}{F}$, $\iotaPlus{\vartheta}{z}{v}{F}$}{descent by one step; iterated descent}
\nt{$\descMove{y}{F}$, $\liftMove{y}{F}$}{descent and lift with the indicator of~$y$}
\nt{$\Sigma$}{a finite set of finite places}
\nt{$V$, $Z$}{the two grids ($z=0$, $z=1$) in Section~\ref{s:grid}; $V$ also a standard space in Theorem~C and in Section~\ref{s:branching}}
\nt{$\Delta_j$, $\mathbf{k}$, $I$, $\mathcal{S}$, $\mathcal{T}$}{a bijection; an element of $\Kmax$; an index set of a
family; spaces}
\nt{$\imag$, $\pi$, $e$}{a square root of $-1$; an integer of norm $\levelN{D}$ (Definition~\ref{d:pi}), also the number $\pi$ in Definition~\ref{d:target} and in the proof of Theorem~A(i); Euler's number}
\end{multicols}
\endgroup

\head{References}
{\fontsize{9.4}{12.6}\selectfont
The references are those of the catalogue of quoted results of this write-up, with the state it records.

[W] A.~Weil, \emph{On a certain type of characters of the id\`ele-class group of an algebraic number-field},
Proceedings of the international symposium on algebraic number theory, Tokyo--Nikko 1955 (1956), pp.~1--7,
especially pp.~2--3 on the unit and finite-index extension conditions, equation~(3).

[R] A.~Raghuram, \emph{Notes on the arithmetic of Hecke $L$-functions}, arXiv:2207.03399, §3.1, equation~(12) and
Proposition~13, pp.~9--10, and §4.3.2, Proposition~25, p.~17. (For [W] and [R] the second reading of the
catalogue finds a different domain and compatibility condition than in Input~(I6): relation unresolved.)

[I] Ichino, \emph{Theta lifting for discrete series representations of real unitary groups}, arXiv:2002.09148,
Theorem~4.1 and Lemma~7.8, pp.~35--36. (Cited by the catalogue for the edge statement behind Input~(I16); its
second reading finds that the formal two-sided statement goes beyond the published joint-type criterion, so the
relation is unresolved.)

For Inputs (I11) and (I12) the catalogue has no confirmed reference for the statements used here: reference to be
confirmed.\par}

\end{document}
